% SPDX-License-Identifier: CC-BY-4.0
% Copyright 2025 Florian Aram Feuerriegel — kassensturz.org
%
% Modelldokumentation der Rechenengine (js/rechner/, js/data.js)
% Übersetzen:  cd docs/modell && node messung.js && latexmk -pdf modell.tex
% Alle Messwerte stammen aus generiert/ und werden von messung.js erzeugt.
\documentclass[11pt,a4paper]{article}

\usepackage[T1]{fontenc}
\usepackage[utf8]{inputenc}
\usepackage{lmodern}
\usepackage{textcomp}
\usepackage{newunicodechar}
% Zeichen aus den Quellenobjekten der Engine (Anhang), für pdfLaTeX übersetzt
\newunicodechar{·}{\textperiodcentered}
\newunicodechar{×}{\ensuremath{\times}}
\newunicodechar{−}{\ensuremath{-}}
\newunicodechar{Δ}{\ensuremath{\Delta}}
\newunicodechar{ε}{\ensuremath{\varepsilon}}
\newunicodechar{α}{\ensuremath{\alpha}}
\newunicodechar{Σ}{\ensuremath{\Sigma}}
\newunicodechar{∑}{\ensuremath{\sum}}
\newunicodechar{∫}{\ensuremath{\int}}
\newunicodechar{÷}{\ensuremath{\div}}
\newunicodechar{½}{\ensuremath{\tfrac12}}
\newunicodechar{²}{\ensuremath{^2}}
\newunicodechar{₀}{\ensuremath{_0}}
\newunicodechar{₂}{\ensuremath{_2}}
\newunicodechar{ₘ}{\ensuremath{_m}}
\newunicodechar{ᵢ}{\ensuremath{_i}}
\newunicodechar{ˣ}{\ensuremath{^x}}
\newunicodechar{≤}{\ensuremath{\leq}}
\newunicodechar{≥}{\ensuremath{\geq}}
\newunicodechar{≈}{\ensuremath{\approx}}
\newunicodechar{→}{\ensuremath{\rightarrow}}
\newcommand{\formelschrift}{\ttfamily\small}
\usepackage[ngerman]{babel}
\usepackage[autostyle]{csquotes}
\usepackage{microtype}
\usepackage[margin=2.4cm]{geometry}
\usepackage{amsmath,amssymb,mathtools}
\usepackage{booktabs,longtable,tabularx,array,multirow}
\usepackage{graphicx}
\usepackage{eurosym}
\usepackage{xcolor}
\usepackage{enumitem}
\usepackage{tikz}
\usetikzlibrary{positioning,arrows.meta,calc,fit,backgrounds,shapes.geometric}
\usepackage{pgfplots}
\pgfplotsset{compat=1.18}
\usepackage{caption}
\usepackage[hidelinks]{hyperref}

\setlength{\parindent}{0pt}
\setlength{\parskip}{0.55em}
\setlist{itemsep=0.2em,topsep=0.3em}

% Farben: positiv (Saldo verbessert / Größe steigt) und negativ, farbfehlsichtig unterscheidbar
\definecolor{pos}{HTML}{1F6FB2}
\definecolor{neg}{HTML}{C4501A}
\definecolor{grau}{HTML}{6B6B6B}
\definecolor{hell}{HTML}{EEF2F6}

\newcommand{\COz}{CO\textsubscript{2}}
\newcommand{\Mrd}{Mrd.\,\euro{}}
\newcommand{\datei}[1]{\texttt{#1}}
\newcommand{\Pp}{Pp.}

% Merksätze am Ende eines Abschnitts
\newenvironment{merksatz}{\par\medskip\noindent\begin{tabular}{@{}p{\linewidth}@{}}\toprule\textbf{Kern.}\ }{\\\bottomrule\end{tabular}\par\medskip}

% Erzeugt von docs/modell/messung.js — nicht von Hand bearbeiten
\newcommand{\SQsaldo}{$-$173,0}
\newcommand{\SQgini}{0,312}
\newcommand{\SQarmut}{15,7}
\newcommand{\SQemiss}{649}
\newcommand{\SQbip}{4.470}
\newcommand{\SQschuldEnde}{102,5}
\newcommand{\SQbipEnde}{6.771}
\newcommand{\SQcoEnde}{6.799}
\newcommand{\SQeinnahmen}{1.679}
\newcommand{\SQausgaben}{1.852}
\newcommand{\SQsaldoPct}{$-$3,87}
\newcommand{\SQstruktur}{$-$1,28}
\newcommand{\Wstabil}{0,372}
\newcommand{\Warbeit}{29,1}
\newcommand{\Winvestlang}{15,6}
\newcommand{\Winvestn}{3,9}
\newcommand{\Wanpassung}{0,252}
\newcommand{\Wschuld}{$-$0,089}


\title{Kassensturz Planspiel\\[0.3em]\large Die Rechenengine in Worten, in Formeln und als Wirkungsgraph}
\author{Modelldokumentation zu \datei{js/rechner/} und \datei{js/data.js}}
\date{Stand: 25.\,September 2026}

\begin{document}
\maketitle

\begin{abstract}
\noindent Das Planspiel lässt Studierende Steuer- und Sozialpolitik einstellen und rechnet
sofort die Folgen für Haushaltssaldo, Verteilung, Emissionen und Wirtschaftsleistung. Dieses
Dokument beschreibt die Rechenengine dreimal: \emph{Teil~I} in Worten, für alle, die wissen
wollen, was das Modell tut, ohne Code zu lesen; \emph{Teil~II} in Formeln, als vollständige
Spezifikation mit Quellen; \emph{Teil~III} als gerichteten, gewichteten Graphen, dessen
Kantengewichte nicht behauptet, sondern gemessen sind: Jede Stellgröße wurde im Modell um
einen festen Schritt verschoben, und die Differenz zum Status quo ist das Gewicht. Alle
Messwerte erzeugt \datei{docs/modell/messung.js}; ein Test meldet, wenn das Dokument nicht
mehr zur Engine passt.
\end{abstract}

\tableofcontents
\clearpage

\section*{Zum Gebrauch}
\addcontentsline{toc}{section}{Zum Gebrauch}

\textbf{Für wen.} Für Lehrende, die das Planspiel einsetzen und seine Ergebnisse im
Debriefing erklären müssen, und für Fachleute, die das Modell prüfen. Teil~I setzt kein
Vorwissen voraus, Teil~II volkswirtschaftliches Grundstudium.

\textbf{Was gilt.} Maßgeblich ist der Code. Wo dieses Dokument und die Engine voneinander
abweichen, irrt das Dokument. Die Zahlen in Teil~III und im Anhang sind aus der Engine
erzeugt und damit aktuell, solange \datei{tests/modelldoku.test.js} grün ist.

\textbf{Einheiten.} Geldgrößen des Staates in Mrd.\,\euro{} je Jahr, Geldgrößen der
Haushalte in \euro{} je Haushalt und Jahr, in Preisen von 2025. Sätze in Prozent, Änderungen
von Sätzen in Prozentpunkten (\Pp). Emissionen in Mt \COz-Äquivalenten. Der Gini-Koeffizient
wird in Teil~III in Gini-Punkten angegeben (1 Punkt = 0,01).

\textbf{Status quo.} Alle Wirkungen sind Differenzen gegenüber dem Status quo, dem
Rechtsstand 2026 (\datei{docs/RECHTSSTAND.md}). In der ersten Periode (2025--2028) liefert er
einen Saldo von \SQsaldo\,\Mrd{} je Jahr (\SQsaldoPct\,\% des BIP), einen Gini von \SQgini{},
eine Armutsrisikoquote von \SQarmut\,\% und Emissionen von \SQemiss\,Mt. Bis zur fünften
Periode (2041--2044) steigt die Schuldenquote ohne Politikänderung auf \SQschuldEnde\,\%.

\clearpage
\part{Das Modell in Worten}
% SPDX-License-Identifier: CC-BY-4.0
% Copyright 2025 Florian Aram Feuerriegel — kassensturz.org
% Teil I — Das Modell in Worten

\section{Was für ein Modell das ist}

Die Engine ist ein \emph{Mikrosimulationsmodell mit makroökonomischem Rahmen}. Sie bildet
Deutschland durch zwölf typische Haushalte ab, je einen für die Einkommensdezile D1 bis D9 und
drei für das oberste Dezil (D10a: 90.--95.~Perzentil, D10b: 95.--99., D10c: das oberste
Prozent). Zusammen stehen sie für rund 41~Millionen Haushalte. Für jeden dieser Haushalte
rechnet die Engine Steuern, Beiträge, Transfers und Nettoeinkommen aus; die Summe über alle
Haushalte, gewichtet mit ihrer Zahl, ergibt die Einnahmen und Ausgaben des Staates.

Darüber liegt ein einfacher Makrorahmen: ein nominales Bruttoinlandsprodukt, das mit einem
Trend wächst, ein Schuldenstand, der sich aus dem Saldo fortschreibt, und ein
Emissionspfad. Die Zeit läuft in \emph{Perioden}, im Standardkurs fünf zu je vier Jahren
(2025--2044), entsprechend einer Legislaturperiode. Innerhalb einer Periode gilt eine Politik
unverändert; zwischen den Perioden wird der Zustand fortgeschrieben.

Drei Eigenschaften bestimmen, wie die Ergebnisse zu lesen sind:

\begin{enumerate}
\item \textbf{Alles ist ein Vergleich.} Die Engine sagt nicht, wie hoch das Nettoeinkommen
  eines Haushalts 2030 sein \emph{wird}, sondern wie es sich gegenüber dem Status quo
  \emph{derselben} Periode unterscheidet. Die Absolutwerte sind kalibriert, aber die Aussagen
  des Planspiels liegen in den Differenzen.
\item \textbf{Verhalten ist berücksichtigt, aber einfach.} Haushalte arbeiten etwas weniger,
  wenn ihnen von einem zusätzlichen Euro weniger bleibt; sie konsumieren weniger, wenn die
  Umsatzsteuer steigt; Unternehmen investieren weniger bei höheren Gewinnsteuern; Emissionen
  sinken mit dem \COz-Preis. Jede dieser Reaktionen ist eine einzige Elastizität aus der
  Literatur, keine geschätzte Nutzenfunktion. Es gibt keine Erwartungen, keine Optimierung
  und keine Rückwirkung auf Löhne oder Preise.
\item \textbf{Jede Zahl hat eine Quelle.} Elastizitäten, Faktoren und Tarifformeln stehen in
  den Objekten \datei{FORMEL\_QUELLEN\_*} und \datei{ELAST\_QUELLEN} mit Fundstelle. Der Anhang
  druckt sie vollständig ab. Wo keine Primärquelle erreichbar war, steht \enquote{Näherung}.
\end{enumerate}

\begin{merksatz}
Zwölf Haushaltstypen, ein Staat, ein Makrorahmen, Perioden zu vier Jahren. Ergebnisse sind
Unterschiede zum Status quo derselben Periode.
\end{merksatz}

\section{Eine Periode in zehn Schritten}
\label{sec:zehn}

Was geschieht, wenn ein Team eine Politik beschließt? Die Funktion \datei{berechne()} geht
für jede Periode dieselben zehn Schritte, Abbildung~\ref{fig:ablauf} zeigt sie im Überblick.

\begin{enumerate}[label=\arabic*.,leftmargin=2em]
\item \textbf{Geltendes Recht einrechnen.} Zu den Reglern kommen die beschlossenen Pfade
  hinzu, gemittelt über die Jahre der Periode: die Senkung der Körperschaftsteuer um je einen
  Punkt ab 2028 und der steigende Rentenbeitrag nach §\,158 SGB~VI. Die Regler zeigen die Werte
  von 2026.
\item \textbf{Arbeitsangebot.} Für jeden Haushalt wird der Grenzsteuersatz mit dem des Status
  quo verglichen. Bleibt vom letzten Euro mehr übrig, steigt das Arbeitseinkommen etwas.
\item \textbf{Einkommensteuer.} Auf das angepasste Arbeitseinkommen wird der Tarif des
  §\,32a EStG angewandt, auf Kapitaleinkommen die Abgeltungsteuer.
\item \textbf{Weitere Steuern.} Körperschaft- und Gewerbesteuer auf die Unternehmensgewinne,
  Umsatzsteuer auf den Konsum, Erbschaft-, Vermögen- und Bodensteuer, \COz-Abgabe auf die
  bepreisten Emissionen.
\item \textbf{Sozialbeiträge} auf die Löhne bis zur Beitragsbemessungsgrenze.
\item \textbf{Transfers und Renten:} Bürgergeld, Kindergeld, gegebenenfalls Klimageld; die
  Renten folgen dem Rentenniveau.
\item \textbf{Haushalte.} Aus allem entsteht das Nettoeinkommen jedes Haushaltstyps, dazu
  die Last aus Unternehmens- und Vermögensteuern, die am Ende Menschen tragen. Daraus folgen
  Gini, Palma und Armutsrisikoquote.
\item \textbf{Nachfrage.} Haushalte, die mehr (weniger) Geld haben als im Status quo, geben
  einen Teil davon aus (weniger aus). Zusammen mit zusätzlichen öffentlichen Investitionen
  verschiebt das die Wirtschaftsleistung dieser Periode.
\item \textbf{Fortschreibung in Euro der Periode.} Alles bis hierher ist in Größen von 2025
  gerechnet. Die Einnahmen werden mit dem tatsächlichen BIP der Periode hochgerechnet, die
  Ausgaben mit dem Preis- und Lohntrend. Dazu kommen Zinsen, Verteidigung und
  Sondervermögen.
\item \textbf{Saldo und Schuldenbremse.} Einnahmen minus Ausgaben ergeben den Saldo, bereinigt
  um Konjunktur und Ausnahmen den strukturellen Saldo, der gegen die Schuldenbremse geprüft
  wird.
\end{enumerate}

Danach übergibt \datei{berechneTransition()} an die nächste Periode: Das BIP wächst mit dem
Trend und den Angebotseffekten, die Schulden wachsen mit Zins und Defizit, die Emissionen
summieren sich zum verbrauchten \COz-Budget.

\begin{figure}[t]
\centering
\begin{tikzpicture}[
  node distance=3.5mm and 5mm,
  box/.style={draw=grau, rounded corners=2pt, fill=hell, align=center, font=\small,
              minimum height=8mm, text width=27mm},
  pfeil/.style={-{Stealth[length=2mm]}, thick, grau}]
\node[box] (a) {1 Geltendes\\Recht};
\node[box, right=of a] (b) {2 Arbeits-\\angebot};
\node[box, right=of b] (c) {3 Einkommen-\\steuer};
\node[box, right=of c] (d) {4 Weitere\\Steuern};
\node[box, below=of d] (e) {5 Sozial-\\beiträge};
\node[box, left=of e] (f) {6 Transfers\\und Renten};
\node[box, left=of f] (g) {7 Haushalte:\\Netto, Gini};
\node[box, left=of g] (h) {8 Nachfrage\\der Periode};
\node[box, below=of h] (i) {9 Fortschreibung\\nominal};
\node[box, right=of i] (j) {10 Saldo,\\Schuldenbremse};
\node[box, right=of j, fill=white, text width=58mm] (k) {Übergang zur nächsten Periode:\\BIP, Schulden, Emissionen, Kapital};
\draw[pfeil] (a)--(b); \draw[pfeil] (b)--(c); \draw[pfeil] (c)--(d); \draw[pfeil] (d)--(e);
\draw[pfeil] (e)--(f); \draw[pfeil] (f)--(g); \draw[pfeil] (g)--(h); \draw[pfeil] (h)--(i);
\draw[pfeil] (i)--(j); \draw[pfeil] (j)--(k);
\end{tikzpicture}
\caption{Ablauf einer Periode in \datei{berechne()} und Übergang in \datei{berechneTransition()}.}
\label{fig:ablauf}
\end{figure}

\section{Die Haushalte}

Jeder der zwölf Haushaltstypen ist durch wenige Zahlen beschrieben (Tabelle~\ref{tab:dezile}
im Anhang): Bruttoeinkommen, Anteil der Kapitaleinkommen, Konsumquote, Vermögen, Zahl der
Haushalte, das Bedarfsgewicht nach der neuen OECD-Skala und der Anteil der gesetzlichen Rente
am Einkommen. Grundlage sind SOEP v40, der Mikrozensus 2024, die Verteilungsrechnung DINA-DE
und der Vermögensbericht des DIW.

Das oberste Dezil ist dreigeteilt, weil sich dort fast alles abspielt, was Vermögen- und
Spitzensteuern betrifft: D10c, das oberste Prozent, hat im Mittel 700.000\,\euro{}
Bruttoeinkommen, zu 45\,\% aus Kapital, und 7~Mio.\,\euro{} Vermögen.

Zu jedem Haushaltstyp gibt es außerdem \emph{Profile}, die Haushalte und Staat gemeinsam
benutzen: Kinder mit Kindergeldanspruch (zusammen 17~Mio.), Bürgergeldbezug (zusammen
3,9~Mio. Regelsatz-Jahresäquivalente, nur D1 bis D4), die Last aus dem \COz-Preis
(regressiv: 4\,\% des Einkommens in D1, 1\,\% in D10c) und die Grenzkonsumneigung. Dass
beide Seiten dieselben Zahlen benutzen, ist nicht selbstverständlich: Bis September 2026
zählten die Haushalte 36,5~Mio. Kinder, der Staat 17~Mio. Seitdem prüft
\datei{tests/buchung.test.js}, dass die Summe über die Haushalte der Buchung des Staates
entspricht.

\begin{merksatz}
Zwölf Haushaltstypen mit Einkommen, Vermögen, Kindern, Transferbezug, Rente und
Konsumverhalten. Staat und Haushalte rechnen mit denselben Profilen.
\end{merksatz}

\section{Die Einkommensteuer}

Der Tarif ist der amtliche Formeltarif des §\,32a EStG 2026: steuerfrei bis zum
Grundfreibetrag (12.348\,\euro{}), dann eine Progressionszone, in der der Grenzsteuersatz
linear von 14 auf knapp 24\,\% steigt, eine zweite bis 42\,\%, eine Proportionalzone mit 42\,\%
und ab 277.826\,\euro{} der Spitzensatz von 45\,\%.

Vier Regler greifen in den Tarif ein, und jeder wirkt nur dort, wo sein Name hinzeigt: Der
\emph{Grundfreibetrag} verschiebt den ganzen Tarif, der \emph{Eingangssatz} hebt oder senkt den
Beginn der Progression, der \emph{Spitzensatz} gilt nur oberhalb der \emph{Grenze}, und die
Grenze verschiebt nur den Beginn der obersten Zone. Das klingt selbstverständlich, war es
aber nicht: Bis September 2026 hob ein höherer Spitzensatz die 42-\%-Zone mit an, und ein
Haushalt aus D5 zahlte für eine Reform der Spitzenverdiener mit.

Weil die Daten \emph{Haushalts}einkommen sind, der Tarif aber je Steuerpflichtigem gilt,
rechnet die Engine um: Nur 79\,\% des Bruttos sind zu versteuerndes Einkommen (Werbungskosten,
Vorsorgeaufwand), und ein Haushalt besteht im Mittel aus 1,6 Tarifeinheiten (Einzelpersonen
und zusammenveranlagte Paare). So trifft das Aufkommen die amtliche Größenordnung.

Für das oberste Prozent reicht ein Durchschnitt nicht: Sein mittleres zu versteuerndes
Einkommen liegt unter der Grenze des Spitzensatzes. Rechnete man nur mit dem Durchschnitt,
zahlte niemand im Modell den Spitzensatz. Die Engine nimmt deshalb an, dass die Einkommen
innerhalb von D10c einer Pareto-Verteilung folgen (Exponent 1,5, der schwere Rand der
Schätzungen für Deutschland), und berechnet, welcher Teil davon über der Grenze liegt.

\begin{merksatz}
Tarif 2026 nach §\,32a EStG, umgerechnet auf Haushalte. Jeder Regler wirkt nur auf seine
Zone. Das oberste Prozent ist als Verteilung modelliert, nicht als Durchschnitt.
\end{merksatz}

\section{Verhalten}

Vier Reaktionen machen aus der mechanischen Rechnung eine Verhaltensrechnung.

\textbf{Arbeitsangebot.} Steigt der Anteil, der vom letzten verdienten Euro übrig bleibt
(eins minus Grenzsteuersatz), um 10\,\%, steigt das Arbeitseinkommen um 2\,\%
(Elastizität 0,20, Konsens von Saez, Chetty und Gruber). Für das oberste Prozent ist die
Elastizität doppelt so hoch (0,40), weil dort Einkommen leichter verlagert werden kann;
oberhalb eines Grenzsteuersatzes von 45\,\% kommt Steuervermeidung hinzu, oberhalb von 60\,\%
Wegzug. Alle Reaktionen sind nach unten und oben begrenzt.

\textbf{Konsum.} Eine um einen Punkt höhere Umsatzsteuer senkt den besteuerten Konsum um
0,35\,\%, getrennt für den Regel- und den ermäßigten Satz.

\textbf{Investitionen.} Liegt die Summe aus Körperschaft- und Gewerbesteuersatz um zehn Punkte
über 30\,\%, sinken Gewinne und Investitionen um 4\,\% (Elastizität $-0{,}40$, Meta-Analyse
von Gechert und Heimberger). Das mindert die Bemessungsgrundlage sofort und das private
Kapital langfristig.

\textbf{Emissionen und Vermögen.} Ein um 100\,\euro{}/t höherer \COz-Preis senkt die bepreisten
Emissionen um 30\,\%. Eine Vermögensteuer senkt das angegebene Vermögen, nach Erfahrungen aus
der Schweiz (Brülhart u.\,a. 2022), hier halbiert.

\begin{merksatz}
Arbeit, Konsum, Investitionen, Emissionen und Vermögen reagieren über je eine Elastizität aus
der Literatur. Alle Reaktionen sind begrenzt, damit extreme Regler nicht ins Unsinnige
führen.
\end{merksatz}

\section{Die übrigen Einnahmen}

\textbf{Umsatzsteuer.} 70\,\% des Konsums fallen unter den Regelsatz, 30\,\% unter den
ermäßigten. Weil der Konsum der Haushalte nur einen Teil der Bemessungsgrundlage erfasst
(dazu kommen Käufe von Staat, Wohnungsbau und befreiten Branchen), wird hochgerechnet, und
3,7\,\% gehen durch Hinterziehung verloren (VAT Gap der EU-Kommission).

\textbf{Unternehmenssteuern.} Körperschaftsteuer (15\,\%) und Gewerbesteuer (effektiv 14\,\%)
haben getrennte Bemessungsgrundlagen, weil Personengesellschaften Gewerbe-, aber keine
Körperschaftsteuer zahlen. Beide schrumpfen mit der Investitionsreaktion.

\textbf{Vermögensbezogene Steuern.} Die Erbschaftsteuer rechnet mit 400\,\Mrd{} Erbmasse im
Jahr, von der nach den persönlichen Freibeträgen 45\,\% steuerpflichtig bleiben;
Betriebsvermögen ist weitgehend verschont. Eine Vermögensteuer trifft Vermögen über
2~Mio.\,\euro{} (3.500\,\Mrd), eine Bodenwertsteuer den Bodenwert (5.000\,\Mrd). Die
Mindeststeuer auf Milliardäre nach Zucman rechnet die gezahlte Einkommensteuer (rund 0,3\,\%
des Vermögens) und eine gleichzeitige Vermögensteuer an.

\textbf{\COz-Preis.} Die Emissionen folgen einem Basispfad nach dem Projektionsbericht 2025
des Umweltbundesamts: 649~Mt im Jahr 2025, $-63\,\%$ bis 2030 und $-80\,\%$ bis 2040 gegenüber
1990. Nur rund die Hälfte davon (327~Mt) trägt den nationalen \COz-Preis von 55\,\euro{}/t;
nur auf diesen Teil wirkt der Regler. Das Aufkommen fließt in den Klimafonds; ein Klimageld
gibt es im Status quo nicht, als Schalter aber schon: Es zahlt 70\,\% des Aufkommens pro
Haushalt gleich aus.

\textbf{Sozialbeiträge.} Renten-, Kranken-, Arbeitslosen- und Pflegeversicherung auf die
sozialversicherungspflichtige Lohnsumme von 1.750\,\Mrd, bis zur Beitragsbemessungsgrenze.
Der Beitragssatz wirkt \emph{nur auf die Einnahmen}. Das ist eine bewusste Entscheidung: Bis
September 2026 folgten auch die Ausgaben dem Satz, und eine Senkung des Rentenbeitrags
verbesserte den Saldo und machte alle Haushalte reicher.

\begin{merksatz}
Jede Steuerart hat ihre eigene Bemessungsgrundlage, kalibriert auf das Ist-Aufkommen.
Beitragssätze bestimmen Einnahmen, nicht Leistungen.
\end{merksatz}

\section{Ausgaben, Transfers und Rente}

Die Ausgaben des Staates beginnen mit dem Stand von 2025, grob nach Aufgaben gegliedert
(soziale Sicherung, Gesundheit, Bildung, Verteidigung, Infrastruktur, Verwaltung, Zinsen).
Die Stellgrößen ändern daran:

\begin{itemize}
\item \textbf{Bürgergeld und Kindergeld:} Regelsatz beziehungsweise Betrag je Kind, mal
  Zahl der Berechtigten aus den Haushaltsprofilen, mal zwölf Monate. Unterkunft und
  Mehrbedarfe (10,8\,\Mrd) bleiben fest.
\item \textbf{Rentenniveau:} Die Rentenzahlungen (rund 362\,\Mrd) steigen oder fallen mit dem
  Sicherungsniveau (Status quo 48\,\%). Jede Änderung kommt in genau der Höhe bei den
  Haushalten an, die der Staat bucht, verteilt nach dem Rentenanteil ihres Einkommens.
  Die Alterung erhöht die Rentenausgaben von allein, bis 2041 um ein Viertel.
\item \textbf{Klimageld}, falls eingeschaltet.
\item \textbf{Erhebungskosten:} Ändert sich das Steuersystem, ändern sich die Kosten, es zu
  erheben (etwa 2,5\,\% des Einkommensteueraufkommens, 8\,\% der Erbschaftsteuer). In den
  Saldo geht nur die Änderung gegenüber dem Status quo.
\end{itemize}

Dazu kommt geltendes Recht, das über die Jahre wirkt: Die Verteidigungsausgaben steigen nach
dem Plan der Bundesregierung von 2,4 auf 3,5\,\% des BIP (2029), das Sondervermögen
Infrastruktur gibt von 2026 bis 2036 jährlich 45\,\Mrd{} aus, und der Klimafonds schrumpft
mit den Emissionen, aus deren Bepreisung er sich speist.

\begin{merksatz}
Ausgaben folgen Leistungen, nicht Beiträgen. Was der Staat für Haushalte ausgibt, kommt dort
in derselben Höhe an.
\end{merksatz}

\section{Vom Staat zum Haushalt: wer trägt was}

Das Nettoeinkommen eines Haushalts ist
\begin{center}
Brutto $-$ Einkommensteuer $-$ Beiträge $-$ Umsatzsteuer $-$ \COz-Last $+$ Transfers $+$
Rentenänderung $-$ Last aus Unternehmens- und Vermögensteuern.
\end{center}

Zwei Posten verdienen eine Erklärung. \emph{Wer trägt eine Beitragserhöhung?} Formal zahlen
Arbeitnehmer und Arbeitgeber je die Hälfte. Langfristig wälzen Arbeitgeber ihren Anteil
aber über niedrigere Löhne auf die Beschäftigten ab (Meta-Analyse von Melguizo und
González-Páramo). Die Engine lässt deshalb eine \emph{Änderung} der Beitragssätze voll bei
den Haushalten ankommen.

\emph{Wer trägt die Körperschaftsteuer?} Unternehmen sind keine Menschen; ihre Steuern tragen
am Ende Beschäftigte über niedrigere Löhne oder Eigentümer über niedrigere Gewinne. Nach
Fuest, Peichl und Siegloch (2018) tragen die Beschäftigten rund die Hälfte der
Gewerbesteuer. Die Engine verteilt eine Änderung von Körperschaft- und Gewerbesteuer je zur
Hälfte nach Arbeits- und Kapitaleinkommen, Erbschaft-, Vermögen- und Bodensteuer nach dem
Vermögen, die Milliardärssteuer auf D10c.

Aus den zwölf Nettoeinkommen entstehen die Verteilungsmaße. Wie in der amtlichen Statistik
(EU-SILC) werden sie auf das \emph{Äquivalenzeinkommen} bezogen: Ein Haushalt mit vier
Personen braucht mehr als einer mit einer, also wird durch das Bedarfsgewicht geteilt.
Der \emph{Gini-Koeffizient} misst die Fläche zwischen Lorenzkurve und Gleichverteilung (0 =
alle gleich, 1 = einer hat alles), die \emph{Palma-Ratio} teilt das Einkommen der oberen
10\,\% durch das der unteren 40\,\%. Die \emph{Armutsrisikoquote} zählt, wer weniger als 60\,\%
des mittleren Einkommens hat; weil ein Durchschnitt je Dezil dafür zu grob ist, wird der
Anteil innerhalb der unteren drei Dezile aus dem Status quo fortgeschrieben.

\begin{merksatz}
Jede Steuer und jede Leistung landet bei einem Haushaltstyp. Verteilungsmaße werden wie in
der amtlichen Statistik auf Äquivalenzeinkommen gerechnet.
\end{merksatz}

\section{Nachfrage}

Wenn Haushalte mehr Geld haben, geben sie einen Teil davon aus, und dieser Teil ist unten
größer als oben: Ein Haushalt in D1 gibt von einem zusätzlichen Euro 65~Cent aus, einer in
D10c 15~Cent (Grenzkonsumneigung nach Jappelli und Pistaferri). Die Engine summiert die
Änderung der Nettoeinkommen gegenüber dem Status quo, gewichtet mit dieser Neigung, und
rechnet sie mit dem Steuer- und Transfermultiplikator (0,6 bei mittlerer Neigung, Gechert
2015) in eine Nachfrageänderung um. Zusätzliche öffentliche Investitionen wirken mit dem
Multiplikator 1,0.

Die Nachfrageänderung wirkt \emph{nur in der laufenden Periode}: Sie hebt oder senkt das BIP
und über das BIP die Einnahmen (etwa 37~Cent je Euro, der automatische Stabilisator). Sie
wirkt symmetrisch: Wer konsolidiert, zahlt mit Wachstum; wer entlastet, gewinnt welches.
Dauerhaft bleibt davon nur, was in öffentliches Kapital fließt.

\begin{merksatz}
Entlastung unten wirkt stärker auf die Nachfrage als oben. Nachfrage wirkt eine Periode,
Kapital bleibt.
\end{merksatz}

\section{Von Periode zu Periode}

\textbf{Wirtschaftsleistung.} Das nominale BIP wächst im Trend um 2,5\,\% im Jahr (2\,\%
Inflation, 0,5\,\% real). Dazu kommen drei Niveaueffekte der Angebotsseite:
\begin{itemize}
\item \emph{Arbeit}: Mehr Arbeitsangebot hebt das BIP sofort, mit dem Arbeitsanteil 0,65.
\item \emph{Privates Kapital}: Höhere Gewinnsteuern senken das Kapital langfristig (mit dem
  Kapitalanteil 0,35), aber langsam: Der Kapitalstock passt sich um 7\,\% im Jahr an.
\item \emph{Öffentliches Kapital}: Zusätzliche öffentliche Investitionen, auch das
  Sondervermögen, bauen einen Kapitalstock auf, der mit 4\,\% im Jahr abschreibt. Ein
  zusätzlicher Stock in Höhe des heutigen öffentlichen Kapitals (1.600\,\Mrd) hebt das BIP um
  8\,\% (Bom und Ligthart 2014).
\end{itemize}
Alle drei sind Niveaus, keine Wachstumsraten: Eine höhere Körperschaftsteuer senkt das BIP
auf ein niedrigeres Niveau, aber nicht jede Periode weiter.

\textbf{Schulden.} Die Schuldenstandsquote wächst jedes Jahr mit dem Zins von 2,0\,\% und sinkt
mit jedem Primärüberschuss; das BIP im Nenner wächst mit 2,5\,\%. Weil der Zins unter der
Wachstumsrate liegt, stabilisiert sich eine Schuldenquote schon bei einem kleinen
Primärdefizit. Ob \enquote{Sparen alternativlos} ist, entscheidet das Spiel also nicht vorab.

\textbf{Emissionen.} Die jährlichen Emissionen folgen dem Basispfad, verringert um die
Wirkung der Politik, und summieren sich Jahr für Jahr. Dem steht ein deutsches Restbudget von
5.380~Mt gegenüber (Bevölkerungsanteil am globalen Budget für 1,7\,°C nach Forster u.\,a.
2026). Einen Klimaschaden auf das deutsche BIP rechnet das Modell nicht: Deutschland
verursacht 1,5\,\% der Weltemissionen, der Rückkanal auf die eigene Wirtschaft ist
vernachlässigbar. Das Klimaziel steht über Emissionen und Budget im Spiel.

\textbf{Demografie und Schocks.} Ein Rentenlastfaktor nach der 14. Bevölkerungsvorausberechnung
erhöht die Rentenausgaben (1,00 in 2025, 1,25 in 2041). Die Lehrperson kann Schocks setzen:
einen BIP-Einbruch, einen Sprung der Schuldenquote oder einen Zinsanstieg.

\begin{merksatz}
Das BIP wächst mit Trend, Arbeit, privatem und öffentlichem Kapital, alles als Niveaus. Die
Schulden wachsen mit Zins und Defizit. Die Emissionen summieren sich gegen ein Budget.
\end{merksatz}

\section{Schuldenbremse und Tragfähigkeit}

Seit der Reform von 2025 erlaubt die Schuldenbremse Bund und Ländern zusammen ein
strukturelles Defizit von 0,70\,\% des BIP. Strukturell heißt: ohne Konjunktureinfluss (ein
Prozent Nachfragelücke bewegt den Saldo um 0,5\,\% des BIP, Budget-Semielastizität der
EU-Kommission) und ohne die Ausnahmen (Verteidigung über 1\,\% des BIP, Sondervermögen). Im
Status quo der ersten Periode liegt der strukturelle Saldo bei \SQstruktur\,\% des BIP; die
Schuldenbremse ist damit verletzt. Das Modell rechnet den Gesamtstaat, nicht Bund und Länder
getrennt.

Dazu kommen Tragfähigkeitsindikatoren aus \datei{abgeleitet.js}: der Primärsaldo, der
Primärüberschuss, der die Schuldenquote gerade stabilisiert (Domar-Bedingung), die Lücke
zwischen beiden und ein Index der Generationengerechtigkeit aus Schulden und verbrauchtem
\COz-Budget.

\section{Der Status quo als Referenz}

Tabelle~\ref{tab:verlauf} zeigt, wohin der Status quo ohne jede Politikänderung führt. Er ist
nicht stabil: Demografie, Verteidigung und Zinsen treiben das Defizit, und die
Schuldenquote steigt von 63,5\,\% auf \SQschuldEnde\,\%. Die Teams entscheiden also nicht
zwischen \enquote{weiter so} und Veränderung, sondern darüber, wer die Anpassung trägt.

\begin{table}[!htb]
\centering\small
% Erzeugt von docs/modell/messung.js — nicht von Hand bearbeiten
\begin{tabular}{lrrrrrrrr}
\toprule
Periode & BIP & Schulden & Saldo & Saldo & strukt. & Emiss. & CO\textsubscript{2} kum. & Gini \\
 & Mrd.\,\euro{} & \% BIP & Mrd.\,\euro{} & \% BIP & \% BIP & Mt/a & Mt & \\
\midrule
2025-$-$1028 & 4.470 & 63,5 & $-$173,0 & $-$1,87 & $-$1,28 & 649 & 0 & 0,312 \\
2029-$-$1032 & 4.968 & 71,5 & $-$140,7 & $-$1,85 & $-$1,43 & 500 & 2.373 & 0,312 \\
2033-$-$1036 & 5.529 & 82,2 & $-$189,4 & $-$1,23 & $-$1,91 & 399 & 4.199 & 0,312 \\
2037-$-$1040 & 6.143 & 93,4 & $-$191,5 & $-$1,75 & $-$1,25 & 314 & 5.669 & 0,312 \\
2041-$-$1044 & 6.771 & 102,5 & $-$165,4 & $-$1,40 & $-$1,90 & 238 & 6.799 & 0,312 \\
\bottomrule
\end{tabular}

\caption{Status quo über fünf Perioden. BIP und Schuldenquote zu Beginn der Periode, übrige
Größen als Jahreswert der Periode, \COz{} kumuliert seit 2025 zu Periodenbeginn.}
\label{tab:verlauf}
\end{table}

\section{Was das Modell nicht kann}
\label{sec:grenzen}

Ein Lehrmodell muss vereinfachen. Diese Vereinfachungen sollte kennen, wer die Ergebnisse
erklärt:

\begin{itemize}
\item \textbf{Kein allgemeines Gleichgewicht.} Löhne, Preise und Zinsen reagieren nicht auf
  die Politik. Der Zins ist fest (außer im Zinsschock).
\item \textbf{Zwölf Haushaltstypen statt Millionen Haushalten.} Innerhalb eines Dezils sind
  alle gleich, außer im obersten Prozent (Pareto-Rand) und bei der Armutsquote.
\item \textbf{Näherungen} ohne Primärtabelle, weil die Quellen beim Anlegen nicht erreichbar
  waren: Rentenanteile je Dezil, Profil der Grenzkonsumneigung, öffentlicher Kapitalstock,
  Milliardärsvermögen, Pareto-Exponent. Sie sind im Code als solche markiert.
\item \textbf{Restabweichung Haushalt gegen Staat.} Bei Umsatzsteuer und Beiträgen weicht die
  Summe über die Haushalte noch um bis zu 9\,\% von der Buchung des Staates ab.
\item \textbf{Nicht gerechnete Schockwirkungen.} Zwei der fünf Schockeffekte
  (Investitionsrückgang, Emissionsminderung) stehen in der Bibliothek, wirken aber nicht; die
  Oberfläche kennzeichnet sie.
\item \textbf{Schuldenbremse} für den Gesamtstaat, nicht getrennt nach Bund und Ländern;
  EU-Fiskalregeln fehlen.
\end{itemize}

Die Messung für Teil~III hat drei weitere Eigenschaften sichtbar gemacht:

\begin{itemize}
\item \textbf{Die Beitragsbemessungsgrenze wirkt asymmetrisch.} Eine Anhebung um 10.000\,\euro{}
  verbessert den Saldo der ersten Periode um rund 6\,\Mrd, eine Senkung um denselben Betrag
  \emph{verbessert} ihn ebenfalls, um rund 3,5\,\Mrd. Grund: Beim Staat sind Mehreinnahmen
  aus der Grenze nach unten bei null abgeschnitten, bei den Haushalten sinken die Beiträge.
  Die entlasteten Haushalte fragen mehr nach, und der Stabilisator hebt den Saldo. Das ist
  ein Buchungsfehler; er ist zur Behebung vorgemerkt.
\item \textbf{Der Spitzensteuersatz bringt wenig:} rund \wert{zerl}{spitze}{gesamt}\,\Mrd{} je
  Prozentpunkt in der ersten Periode, weniger als übliche Schätzungen (etwa 1 bis 2\,\Mrd).
  Die Zerlegung in Abschnitt~\ref{sec:zerlegung} zeigt, dass das fast ganz an der Rechnung
  selbst liegt (\wert{zerl}{spitze}{mechanisch}\,\Mrd{} mechanisch), nicht am Verhalten
  (\wert{zerl}{spitze}{verhalten}\,\Mrd): Der Satz trifft nur den Teil der Einkommen im
  obersten Prozent, der nach dem Pareto-Rand über der Grenze liegt. Ob der Rand zu dünn ist,
  wäre mit der Einkommensteuerstatistik zu prüfen.
\item \textbf{Die Mindeststeuer auf Milliardäre greift erst oberhalb von 0,3\,\%}, weil die
  gezahlte Einkommensteuer angerechnet wird. Das ist gewollt, erklärt aber, warum kleine
  Sätze nichts bewirken.
\end{itemize}


\clearpage
\part{Das Modell in Formeln}
% SPDX-License-Identifier: CC-BY-4.0
% Copyright 2025 Florian Aram Feuerriegel — kassensturz.org
% Teil II — Das Modell in Formeln
\emph{In Arbeit (Zwischenziel 3).}


\clearpage
\part{Das Modell als Wirkungsgraph}
% SPDX-License-Identifier: CC-BY-4.0
% Copyright 2025 Florian Aram Feuerriegel — kassensturz.org
% Teil III — Das Modell als Wirkungsgraph
%
% Alle Zahlen in Kanten, Tabellen und Kurven stammen aus generiert/ (messung.js).

\section{Wie die Graphen entstehen}
\label{sec:graphmethode}

Teil~II beschreibt das Modell als Gleichungssystem. Man kann dasselbe System als
\emph{gerichteten, gewichteten Graphen} lesen: Knoten sind Größen, eine Kante $u\to v$ heißt
\enquote{$u$ geht in die Berechnung von $v$ ein}, und ihr Gewicht sagt, um wie viel sich $v$
ändert, wenn $u$ sich um eine Einheit ändert. Der Graph hat drei Ebenen:

\begin{enumerate}
\item \textbf{Stellgrößen} (die Regler am Verhandlungstisch),
\item \textbf{Kanäle} (Zwischengrößen: Aufkommen je Steuerart, Transfers, Renten,
  Arbeitsangebot, Investitionen, Nachfrage),
\item \textbf{Kennzahlen} (was das Spiel bewertet: Saldo, Verteilung, Emissionen, BIP,
  Schulden).
\end{enumerate}

Die Gewichte stammen aus zwei Quellen, und der Unterschied ist wichtig:

\begin{itemize}
\item \textbf{Stellgröße $\to$ Kanal: gemessen.} Jede Stellgröße wurde, ausgehend vom Status
  quo, um einen festen Schritt verschoben (Tabelle~\ref{tab:schritte}), die Engine rechnete
  die erste Periode (2025--2028) neu, und die Differenz jedes Kanals zum Status quo ist das
  Gewicht. Formal ist das die endliche Differenz
  \begin{equation}
    w_{u\to v}=\frac{v(\theta^{SQ}+h\,e_u)-v(\theta^{SQ})}{1}\qquad\text{je Schritt }h,
  \end{equation}
  keine Ableitung: Die Engine hat Knicke (Abschnitt~\ref{sec:eigenschaften}), und die
  Wirkung gilt genau für den angegebenen Schritt. Wo das Modell nichtlinear ist, zeigen es
  die Kurven in Abschnitt~\ref{sec:kurven}.
\item \textbf{Kanal $\to$ Kennzahl: strukturell.} Diese Gewichte folgen direkt aus den
  Gleichungen von Teil~II, ausgewertet im Status quo. Ein Euro zusätzliches Aufkommen
  verbessert den Saldo um einen Euro; ein Euro Nachfrage bringt \Wstabil\,\euro{} Einnahmen.
\end{itemize}

Die Summe über alle Pfade von einer Stellgröße zu einer Kennzahl ergibt deren
Gesamtwirkung. Wegen der Rückkopplung über die Nachfrage ist der Graph nicht kreisfrei; die
Engine löst den Kreis, indem sie die Nachfrage einmal gegen den Status quo derselben Periode
rechnet (keine Iteration bis zum Fixpunkt).

\begin{table}[!htb]
\centering\small
\begin{tabular}{lll}
\toprule
Art der Stellgröße & Schritt $h$ & Beispiele \\
\midrule
Steuer- und Beitragssätze & $+1$\,\Pp & Eingangs-, Spitzen-, Umsatz-, Körperschaftsteuer, RV, KV \\
Rentenniveau & $+1$\,\Pp & 48\,\% $\to$ 49\,\% \\
Grundfreibetrag & $+1.000$\,\euro & 12.348 $\to$ 13.348\,\euro \\
Monatliche Leistungen & $+10$\,\euro/Monat & Bürgergeld, Kindergeld \\
\COz-Preis & $+10$\,\euro/t & 55 $\to$ 65\,\euro/t \\
Schalter & aus $\to$ ein & Klimageld \\
\bottomrule
\end{tabular}
\caption{Schritte der Messung. Die Stellgrößen der klassischen Fläche stehen mit ihren
Schritten in Tabelle~\ref{tab:zerlegung}.}
\label{tab:schritte}
\end{table}

\section{Der Strukturgraph}
\label{sec:struktur}

Abbildung~\ref{fig:struktur} zeigt, wie die Kanäle auf die Kennzahlen wirken, unabhängig
davon, welche Stellgröße sie bewegt hat. Er ist das Gerüst, in das sich die gemessenen
Ressortgraphen einhängen.

\begin{figure}[p]
\centering
\resizebox{\textwidth}{!}{%
\begin{tikzpicture}[
  kn/.style={draw=grau, rounded corners=2pt, fill=hell, align=center, font=\small, text width=27mm, minimum height=10mm},
  zk/.style={draw=black, thick, rounded corners=2pt, fill=white, align=center, font=\small, text width=27mm, minimum height=10mm},
  pf/.style={-{Stealth[length=2.2mm]}, thick},
  rk/.style={-{Stealth[length=2.2mm]}, thick, dashed, grau},
  gl/.style={font=\scriptsize, fill=white, inner sep=1pt, align=center}]
% Spalte A: Kanäle
\node[kn] (aus)  at (0, 1.9)   {Ausgaben an Haushalte\\{\scriptsize Transfers, Renten}};
\node[kn] (ein)  at (0,-0.3)   {Einnahmen\\{\scriptsize ESt, USt, KSt, GewSt,\\Verm., \COz, Beiträge}};
\node[kn] (arb)  at (0,-5.8)   {Arbeitsangebot $\bar\lambda$};
\node[kn] (inv)  at (0,-7.7)   {Investitionsfaktor $I_f$};
\node[kn] (emi)  at (0,-9.8)   {Emissionen $E$};
% Spalte B: Haushalte und Nachfrage
\node[kn] (hh)   at (4.3,-2.4) {Nettoeinkommen\\der Haushalte $\Delta_i$};
\node[kn] (nf)   at (4.3,-4.7) {Nachfrageimpuls $J$};
% Spalte C: Kennzahlen der Periode und Niveau
\node[zk] (sal)  at (8.6, 0.8) {Saldo $S$};
\node[zk] (ver)  at (8.6,-2.4) {Gini, Palma,\\Armutsrisiko\\{\scriptsize auf Äquivalenzeinkommen}};
\node[zk] (bip)  at (8.6,-4.7) {BIP der Periode\\$Y\eta$};
\node[zk] (bipl) at (8.6,-6.9) {BIP-Niveau\\Folgeperiode};
\node[zk] (cob)  at (8.6,-9.8) {\COz-Budget\\(kumuliert)};
% Spalte D: Schulden
\node[zk] (sq)   at (12.9,-2.4) {Schuldenquote\\Folgeperiode};
% Kanten der Periode
\draw[pf] (ein.east) -- node[gl, pos=0.5]{$+1$} (sal.west);
\draw[pf] (aus.east) -- node[gl, pos=0.75]{$-1$} (sal.west);
\draw[pf] (ein.east) -- node[gl, pos=0.5, below left=0pt and 1pt]{$-1$ {\tiny Inzidenz}} (hh.west);
\draw[pf] (aus.east) -- node[gl, pos=0.18]{$+1$} (hh.west);
\draw[pf] (hh.east)  -- (ver.west);
\draw[pf] (hh.south) -- node[gl]{$\frac{0{,}6}{0{,}48}\,m_i$ {\tiny(0,19 bis 0,81)}} (nf.north);
\draw[pf] (nf.east)  -- node[gl]{$+1$} (bip.west);
\draw[pf] (nf.north east) -- node[gl, pos=0.78]{$+\Wstabil$} (sal.south west);
% Kanten über die Periode hinaus
\draw[pf] (sal.east) -- node[gl]{$\Wschuld$\,\Pp\\{\tiny je Mrd./a}} (sq.north west);
\draw[pf] (bipl.east) -- node[gl, pos=0.5]{Nenner} (sq.south);
\draw[pf] (arb.east) -- node[gl, pos=0.5]{$+\Warbeit$ {\tiny Mrd. je \%}} (bipl.west);
\draw[pf] (inv.east) -- node[gl, pos=0.5]{$+\Winvestn$ / $\Winvestlang$ {\tiny Mrd. je \%, 1 Per. / lang}} (bipl.south west);
\draw[pf] (emi.east) -- node[gl]{$\times\,n$ Jahre} (cob.west);
% Rückwirkungen (gestrichelt)
\draw[rk] (bipl.north) -- node[gl, text=grau]{{\tiny Ausgangsniveau}\\{\tiny nächste Periode}} (bip.south);
\draw[rk] (arb.west) to[out=160,in=200,looseness=1.0] node[gl, pos=0.5, text=grau]{Basis} (ein.west);
\draw[rk] (inv.west) to[out=170,in=190,looseness=1.3] node[gl, pos=0.5, text=grau]{Basis $\pi$} (ein.south west);
\draw[rk] (emi.west) to[out=175,in=185,looseness=1.5] node[gl, pos=0.5, text=grau]{\COz-Basis} ($(ein.south west)+(0,0.2)$);
\end{tikzpicture}}
\caption{Strukturgraph: wie Kanäle auf Kennzahlen wirken. Durchgezogene Kanten tragen das
Gewicht aus den Gleichungen von Teil~II im Status quo: \euro{} Kennzahl je \euro{} Kanal,
wenn nicht anders angegeben. Gestrichelte Kanten sind Rückwirkungen auf die
Bemessungsgrundlagen (Verhalten) und in die nächste Periode. $m_i$ ist die
Grenzkonsumneigung (0,15 in D10c bis 0,65 in D1).}
\label{fig:struktur}
\end{figure}

\paragraph{So liest man den Strukturgraphen.} Beginnen Sie bei einem Kanal und folgen Sie
den Pfeilen. Beispiel \emph{ein Euro mehr Rente}: Er verschlechtert den Saldo um einen Euro
(Kante \enquote{Ausgaben $\to$ Saldo}, $-1$) und erhöht das Nettoeinkommen eines Haushalts um
einen Euro ($+1$). Dieser Haushalt gibt davon je nach Dezil 15 bis 65~Cent aus; mit dem
Multiplikator werden daraus 19 bis 81~Cent Nachfrage. Jeder Euro Nachfrage hebt das BIP der
Periode um einen Euro und bringt \Wstabil\,\euro{} Einnahmen zurück (Kante
\enquote{Nachfrage $\to$ Saldo}). Ein Euro Rente an einen Haushalt aus D1 kostet den Staat
also nicht einen Euro, sondern rund $1-0{,}81\cdot\Wstabil\approx0{,}70$\,\euro. Über die
Periode hinaus wirkt nur, was die gestrichelten und die unteren Kanten tragen: Arbeit,
Kapital und die Schulden.

\section{Die Ressortgraphen}
\label{sec:ressortgraphen}

Die folgenden vier Abbildungen zeigen für jedes Ressort am Verhandlungstisch, welche Kanäle
seine Stellgrößen bewegen. Die Kanäle stehen in allen vier Graphen an derselben Stelle,
damit man sie vergleichen kann.

\paragraph{So liest man einen Ressortgraphen.}
\begin{itemize}
\item \textbf{Links} steht die Stellgröße mit ihrem Schritt und ihrer Gesamtwirkung auf den
  Saldo der ersten Periode, \textbf{rechts} die Kanäle mit ihrer Einheit.
\item Die \textbf{Zahl an der Kante} ist die gemessene Änderung des Kanals, in seiner
  Einheit, für genau diesen Schritt.
\item \textbf{Blau} heißt: Der Kanal steigt. \textbf{Orange} heißt: Er sinkt. Ob das gut ist,
  hängt vom Kanal ab: Steigende Transfers sind für den Saldo schlecht, für die Haushalte gut.
\item Die \textbf{Strichstärke} ist proportional zum Betrag, gemessen am größten Wert
  derselben Einheit über alle Ressorts. Dicke Kanten sind also in allen vier Graphen gleich
  groß.
\item Kanten unter 0,3\,\Mrd{} beziehungsweise 0,02\,\% sind weggelassen.
\end{itemize}
Beispiel aus Abbildung~\ref{fig:soziales}: Ein Punkt Rentenbeitrag bringt
\wert{kanal}{rv}{sv}\,\Mrd{} Beiträge (dicke blaue Kante), senkt aber das Aufkommen aus
Umsatzsteuer um \wert{kanal}{rv}{mwst}\,\Mrd{} und aus Einkommensteuer um
\wert{kanal}{rv}{est}\,\Mrd, weil die Haushalte weniger Netto haben und weniger konsumieren.
Der Nachfrageimpuls sinkt um \wert{kanal}{rv}{nachfr}\,\Mrd. Unter dem Strich bleiben
\wert{p1}{rv}{saldo}\,\Mrd.

\newcommand{\ressortgraph}[3]{%
\begin{figure}[p]
\centering
\begin{tikzpicture}
\input{generiert/graph_#1.tex}
\end{tikzpicture}
\caption{#2}
\label{#3}
\end{figure}}

\ressortgraph{Finanzen}{Ressort Finanzen und Steuern: gemessene Wirkung auf die Kanäle,
erste Periode, je Schritt.}{fig:finanzen}
\ressortgraph{Wirtschaft}{Ressort Wirtschaft und Unternehmen. Die Investitionsreaktion wirkt
zusätzlich über das private Kapital auf das BIP der Folgeperioden
(Abbildung~\ref{fig:struktur}).}{fig:wirtschaft}
\ressortgraph{Soziales}{Ressort Arbeit und Soziales. Beiträge wirken nur auf die Einnahmen,
das Rentenniveau nur auf die Ausgaben.}{fig:soziales}
\ressortgraph{Umwelt}{Ressort Umwelt und Klima. Das Klimageld verschiebt das
\COz-Aufkommen vollständig zu den Transfers.}{fig:umwelt}

\paragraph{Was die Ressortgraphen zeigen.}
\begin{itemize}
\item \textbf{Jede Steuer bewegt mehr als ihren eigenen Kanal.} Eine höhere Umsatzsteuer
  senkt auch Einkommensteuer und Beiträge, weil sie das Vornetto nicht ändert, wohl aber das
  Konsumvolumen; eine höhere Beitragsrate senkt Einkommen- und Umsatzsteuer.
\item \textbf{Der Nachfrageimpuls ist fast immer die zweitgrößte Kante.} Er ist der Grund,
  warum keine Konsolidierung umsonst ist.
\item \textbf{Nur zwei Stellgrößen am Tisch bewegen das Arbeitsangebot} messbar:
  Grundfreibetrag und Eingangssatz, weil sie den Grenzsteuersatz der Mitte ändern. Der
  Spitzensatz trifft zu wenige Einkommen, und Beiträge gehen im Modell nicht in den
  Grenzsteuersatz ein (Abschnitt~\ref{sec:zerlegung}).
\end{itemize}

\section{Die Gesamtmatrix: Zielkonflikte auf einen Blick}
\label{sec:matrix}

Tabelle~\ref{tab:gesamt} fasst alle Pfade zusammen: Jede Zeile ist eine Stellgröße, jede
Spalte eine Kennzahl, jede Zelle die gemessene Gesamtwirkung. Die Färbung zeigt diesmal nicht
das Vorzeichen, sondern die \emph{Erwünschtheit}: \textcolor{pos}{\textbf{blau}} heißt, die
Kennzahl entwickelt sich in die erwünschte Richtung (höherer Saldo, niedrigerer Gini,
niedrigere Emissionen, höheres BIP, höheres Nettoeinkommen, niedrigere Schuldenquote),
\textcolor{neg}{\textbf{orange}} das Gegenteil. Die Sättigung folgt dem Betrag, relativ zum
größten Wert der Spalte.

\begin{table}[p]
\centering\scriptsize
\setlength{\tabcolsep}{3pt}
% Erzeugt von docs/modell/messung.js — nicht von Hand bearbeiten
\begin{tabular}{lrrrrrrrrrrr}
\toprule
Stellgröße & \rotatebox{70}{Saldo} & \rotatebox{70}{Gini} & \rotatebox{70}{Armutsrisiko} & \rotatebox{70}{Emissionen} & \rotatebox{70}{BIP} & \rotatebox{70}{Netto D1} & \rotatebox{70}{Netto D5} & \rotatebox{70}{Netto D10c} & \rotatebox{70}{Schuldenquote 2041} & \rotatebox{70}{BIP 2041} & \rotatebox{70}{CO\textsubscript{2} 2025–2040} \\
 & {\tiny Mrd. \euro{}/a} & {\tiny Gini-Pkt.} & {\tiny \%} & {\tiny Mt/a} & {\tiny Mrd. \euro{}} & {\tiny \euro{}/a} & {\tiny \euro{}/a} & {\tiny \euro{}/a} & {\tiny Pp.} & {\tiny Mrd. \euro{}} & {\tiny Mt} \\
\midrule
Grundfreibetrag & \cellcolor{neg!45}$-$6,9 & \cellcolor{pos!35}$-$0,13 & \cellcolor{neg!70}0,29 & \textcolor{grau}{0} & \cellcolor{pos!60}8,5 & \textcolor{grau}{0} & \cellcolor{pos!69}358 & \cellcolor{pos!24}637 & \cellcolor{neg!25}0,85 & \cellcolor{pos!70}28,2 & \textcolor{grau}{0} \\
Eingangssteuersatz & \cellcolor{pos!24}2,6 & \cellcolor{pos!19}$-$0,04 & \cellcolor{pos!29}$-$0,09 & \textcolor{grau}{0} & \cellcolor{neg!28}$-$2,8 & \textcolor{grau}{0} & \cellcolor{neg!28}$-$101 & \cellcolor{neg!18}$-$322 & \cellcolor{pos!20}$-$0,51 & \cellcolor{neg!28}$-$7,9 & \textcolor{grau}{0} \\
Spitzensteuersatz & \cellcolor{pos!13}0,3 & \cellcolor{pos!15}$-$0,02 & \textcolor{grau}{0} & \textcolor{grau}{0} & \cellcolor{neg!12}$-$0,1 & \textcolor{grau}{0} & \textcolor{grau}{0} & \cellcolor{neg!33}$-$1.156 & \cellcolor{pos!13}$-$0,09 & \cellcolor{neg!12}$-$0,2 & \textcolor{grau}{0} \\
Umsatzsteuer & \cellcolor{pos!49}7,7 & \cellcolor{neg!22}0,06 & \cellcolor{pos!19}$-$0,03 & \textcolor{grau}{0} & \cellcolor{neg!43}$-$5,5 & \cellcolor{neg!29}$-$87 & \cellcolor{neg!44}$-$204 & \cellcolor{neg!41}$-$1.555 & \cellcolor{pos!48}$-$2,44 & \cellcolor{neg!29}$-$8,2 & \textcolor{grau}{0} \\
Erbschaftsteuer & \cellcolor{pos!13}0,3 & \cellcolor{pos!13}$-$0,01 & \textcolor{grau}{0} & \textcolor{grau}{0} & \cellcolor{neg!13}$-$0,1 & \textcolor{grau}{0} & \cellcolor{neg!12}$-$2 & \cellcolor{neg!15}$-$177 & \cellcolor{pos!13}$-$0,08 & \cellcolor{neg!12}$-$0,2 & \textcolor{grau}{0} \\
Körperschaftsteuer & \cellcolor{pos!22}2,0 & \cellcolor{pos!18}$-$0,04 & \cellcolor{pos!14}$-$0,01 & \textcolor{grau}{0} & \cellcolor{neg!18}$-$1,0 & \cellcolor{neg!14}$-$8 & \cellcolor{neg!16}$-$27 & \cellcolor{neg!42}$-$1.616 & \cellcolor{pos!17}$-$0,33 & \cellcolor{neg!29}$-$8,1 & \textcolor{grau}{0} \\
Gewerbesteuer & \cellcolor{pos!31}3,9 & \cellcolor{pos!25}$-$0,07 & \cellcolor{pos!16}$-$0,02 & \textcolor{grau}{0} & \cellcolor{neg!23}$-$2,0 & \cellcolor{neg!15}$-$16 & \cellcolor{neg!20}$-$52 & \cellcolor{neg!70}$-$3.128 & \cellcolor{pos!26}$-$0,93 & \cellcolor{neg!32}$-$9,6 & \textcolor{grau}{0} \\
RV-Beitrag & \cellcolor{pos!68}11,7 & \cellcolor{neg!24}0,06 & \cellcolor{pos!33}$-$0,10 & \textcolor{grau}{0} & \cellcolor{neg!70}$-$10,3 & \cellcolor{neg!36}$-$129 & \cellcolor{neg!70}$-$366 & \cellcolor{neg!30}$-$987 & \cellcolor{pos!68}$-$3,76 & \cellcolor{neg!44}$-$15,6 & \textcolor{grau}{0} \\
KV-Beitrag & \cellcolor{pos!70}12,1 & \cellcolor{neg!34}0,12 & \cellcolor{pos!33}$-$0,10 & \textcolor{grau}{0} & \cellcolor{neg!66}$-$9,6 & \cellcolor{neg!36}$-$129 & \cellcolor{neg!70}$-$366 & \cellcolor{neg!25}$-$679 & \cellcolor{pos!70}$-$3,90 & \cellcolor{neg!42}$-$14,5 & \textcolor{grau}{0} \\
Rentenniveau & \cellcolor{neg!40}$-$5,8 & \cellcolor{pos!51}$-$0,22 & \cellcolor{pos!13}$-$0,01 & \textcolor{grau}{0} & \cellcolor{pos!39}4,8 & \cellcolor{pos!35}122 & \cellcolor{pos!50}238 & \cellcolor{pos!14}117 & \cellcolor{neg!42}2,01 & \cellcolor{pos!31}9,0 & \textcolor{grau}{0} \\
Bürgergeld & \cellcolor{neg!14}$-$0,4 & \cellcolor{pos!20}$-$0,04 & \cellcolor{pos!27}$-$0,08 & \textcolor{grau}{0} & \cellcolor{pos!14}0,4 & \cellcolor{pos!26}72 & \textcolor{grau}{0} & \textcolor{grau}{0} & \cellcolor{neg!14}0,11 & \cellcolor{pos!13}0,6 & \textcolor{grau}{0} \\
Kindergeld & \cellcolor{neg!20}$-$1,7 & \cellcolor{pos!24}$-$0,07 & \cellcolor{pos!16}$-$0,02 & \textcolor{grau}{0} & \cellcolor{pos!19}1,3 & \cellcolor{pos!20}45 & \cellcolor{pos!21}59 & \cellcolor{pos!12}11 & \cellcolor{neg!20}0,52 & \cellcolor{pos!16}2,0 & \textcolor{grau}{0} \\
CO\textsubscript{2}-Preis & \cellcolor{pos!22}2,0 & \cellcolor{neg!16}0,02 & \cellcolor{pos!14}$-$0,01 & \cellcolor{pos!70}$-$9,8 & \cellcolor{neg!20}$-$1,4 & \cellcolor{neg!17}$-$24 & \cellcolor{neg!21}$-$55 & \cellcolor{neg!18}$-$303 & \cellcolor{pos!19}$-$0,46 & \cellcolor{neg!14}$-$0,8 & \cellcolor{pos!70}$-$103 \\
Klimageld & \cellcolor{neg!58}$-$9,7 & \cellcolor{pos!70}$-$0,32 & \cellcolor{pos!43}$-$0,15 & \textcolor{grau}{0} & \cellcolor{pos!56}7,8 & \cellcolor{pos!70}307 & \cellcolor{pos!61}307 & \cellcolor{pos!18}307 & \cellcolor{neg!45}2,19 & \cellcolor{pos!21}4,3 & \textcolor{grau}{0} \\
\bottomrule
\end{tabular}

\caption{Gesamtmatrix der Wirkungen. Die ersten acht Spalten gelten für die erste Periode
(Jahreswerte 2025--2028), die letzten drei für das Ende des Pfads nach fünf Perioden mit
unveränderter Politik. \textcolor{pos}{Blau}: erwünschte Richtung, \textcolor{neg}{orange}:
unerwünschte. Grau: keine messbare Wirkung. Schritte wie in Tabelle~\ref{tab:schritte}.}
\label{tab:gesamt}
\end{table}

\paragraph{So liest man die Matrix.} Lesen Sie eine \emph{Zeile}: Sie zeigt den Preis einer
Stellgröße. Keine Zeile ist durchgehend blau, das ist die Eigenschaft (E5) aus
Abschnitt~\ref{sec:eigenschaften} und wird von \datei{tests/dominanz.test.js} erzwungen. Beispiel
\emph{Grundfreibetrag}: Er entlastet D5 und D10c, senkt den Gini und hebt das BIP, kostet aber
\wert{p1}{freibetrag}{saldo}\,\Mrd{} Saldo und erhöht die Schuldenquote 2041 um
\wert{pfad}{freibetrag}{schuld}\,\Pp. Dass er die Armutsrisikoquote \emph{erhöht}
(\wert{p1}{freibetrag}{armut}\,\Pp), ist kein Fehler: D1 liegt schon unter dem Freibetrag und
gewinnt nichts, die Mitte gewinnt, das mittlere Einkommen steigt, und mit ihm die relative
Armutsgrenze.

Lesen Sie eine \emph{Spalte}: Sie zeigt, welche Hebel eine Kennzahl überhaupt bewegen. Die
Emissionen bewegt allein der \COz-Preis; das BIP 2041 bewegen am stärksten Freibetrag,
Beiträge und Rentenniveau, über Nachfrage, Arbeitsangebot und Schulden.

\section{Verteilung über die Dezile}

Tabelle~\ref{tab:dezile-wirkung} zeigt die Änderung des Nettoeinkommens jedes Haushaltstyps
in \euro{} im Jahr. Sie macht die Inzidenz sichtbar, die in Gini und Palma nur verdichtet
erscheint.

\begin{table}[!htb]
\centering\scriptsize
\setlength{\tabcolsep}{2.5pt}
% Erzeugt von docs/modell/messung.js — nicht von Hand bearbeiten
\begin{tabular}{lrrrrrrrrrrrr}
\toprule
 & D1 & D2 & D3 & D4 & D5 & D6 & D7 & D8 & D9 & D10a & D10b & D10c \\
\midrule
Grundfreibetrag & 0 & 0 & $+$642 & $+$341 & $+$358 & $+$381 & $+$409 & $+$447 & $+$515 & $+$613 & $+$626 & $+$637 \\
Eingangssteuersatz & 0 & 0 & $-$36 & $-$71 & $-$101 & $-$137 & $-$176 & $-$223 & $-$284 & $-$327 & $-$317 & $-$322 \\
Spitzensteuersatz & 0 & 0 & 0 & 0 & 0 & 0 & 0 & 0 & 0 & 0 & 0 & $-$1.156 \\
Umsatzsteuer & $-$87 & $-$130 & $-$161 & $-$182 & $-$204 & $-$229 & $-$260 & $-$298 & $-$361 & $-$397 & $-$618 & $-$1.555 \\
Erbschaftsteuer & 0 & 0 & 0 & $-$1 & $-$2 & $-$3 & $-$5 & $-$9 & $-$16 & $-$21 & $-$38 & $-$177 \\
Körperschaftsteuer & $-$8 & $-$12 & $-$17 & $-$21 & $-$27 & $-$32 & $-$41 & $-$54 & $-$78 & $-$108 & $-$276 & $-$1.616 \\
Gewerbesteuer & $-$16 & $-$24 & $-$33 & $-$40 & $-$52 & $-$62 & $-$79 & $-$104 & $-$152 & $-$209 & $-$534 & $-$3.128 \\
RV-Beitrag & $-$129 & $-$194 & $-$248 & $-$304 & $-$366 & $-$440 & $-$527 & $-$649 & $-$842 & $-$974 & $-$979 & $-$987 \\
KV-Beitrag & $-$129 & $-$194 & $-$248 & $-$304 & $-$366 & $-$440 & $-$527 & $-$649 & $-$665 & $-$670 & $-$674 & $-$679 \\
Rentenniveau & $+$122 & $+$202 & $+$236 & $+$243 & $+$238 & $+$218 & $+$182 & $+$152 & $+$133 & $+$109 & $+$115 & $+$117 \\
Bürgergeld & $+$72 & $+$30 & $+$10 & $+$2 & 0 & 0 & 0 & 0 & 0 & 0 & 0 & 0 \\
Kindergeld & $+$45 & $+$61 & $+$67 & $+$64 & $+$59 & $+$53 & $+$47 & $+$42 & $+$36 & $+$28 & $+$20 & $+$11 \\
CO\textsubscript{2}-Preis & $-$24 & $-$35 & $-$42 & $-$49 & $-$55 & $-$62 & $-$70 & $-$78 & $-$90 & $-$97 & $-$143 & $-$303 \\
Klimageld & $+$307 & $+$307 & $+$307 & $+$307 & $+$307 & $+$307 & $+$307 & $+$307 & $+$307 & $+$307 & $+$307 & $+$307 \\
\bottomrule
\end{tabular}

\caption{Änderung des Nettoeinkommens je Haushalt und Jahr in \euro, erste Periode, je
Schritt. D10a bis D10c teilen das oberste Dezil (90.--95., 95.--99., oberstes Prozent).}
\label{tab:dezile-wirkung}
\end{table}

\paragraph{So liest man die Tabelle.} Vergleichen Sie Beträge innerhalb einer Zeile mit den
Einkommen aus Tabelle~\ref{tab:dezile}: Die \COz-Abgabe kostet D1 \wert{p1}{co2}{d1}\,\euro{}
und D10c \wert{p1}{co2}{d10c}\,\euro{}; gemessen am Einkommen (14.000 gegen 700.000\,\euro)
trifft sie D1 aber viermal so stark. Das Klimageld zahlt allen gleich viel und kehrt diese
Last deshalb um.

\section{Wirkung über den ganzen Reglerbereich}
\label{sec:kurven}

Die Gewichte oben gelten für einen Schritt am Status quo. Abbildung~\ref{fig:kurven} zeigt
für fünf Stellgrößen, wie Saldo und Verteilung (beim \COz-Preis: Emissionen) über den
ganzen Bereich des Reglers verlaufen; alle anderen Regler bleiben im Status quo. Die
gestrichelte Linie markiert den Status quo.

\newcommand{\kurve}[6]{% Stellgröße, Spalte, x-Beschriftung, y-Beschriftung, Titel, Farbe
  \nextgroupplot[title={#5}, xlabel={#3}, ylabel={#4},
    extra x ticks={\wert{sq}{#1}{x}}, extra x tick labels={},
    extra x tick style={grid=major, major grid style={dashed, grau}}]
  \addplot[color=#6, very thick, mark=none] table[x=x, y=#2] {generiert/verlauf_#1.dat};}

\begin{figure}[p]
\centering
\begin{tikzpicture}
\begin{groupplot}[group style={group size=2 by 5, horizontal sep=18mm, vertical sep=21mm},
  width=0.44\textwidth, height=0.15\textheight,
  tick label style={font=\scriptsize}, label style={font=\scriptsize},
  title style={font=\small\bfseries}, grid=none,
  /pgf/number format/use comma, /pgf/number format/1000 sep={.},
  scaled y ticks=false]
\kurve{spitze}{saldo}{Spitzensteuersatz [\%]}{Saldo [\Mrd]}{Spitzensteuersatz: Saldo}{pos}
\kurve{spitze}{gini}{Spitzensteuersatz [\%]}{Gini [Pkt.]}{Spitzensteuersatz: Gini}{neg}
\kurve{mwst}{saldo}{Umsatzsteuer [\%]}{Saldo [\Mrd]}{Umsatzsteuer: Saldo}{pos}
\kurve{mwst}{gini}{Umsatzsteuer [\%]}{Gini [Pkt.]}{Umsatzsteuer: Gini}{neg}
\kurve{kst}{saldo}{Körperschaftsteuer [\%]}{Saldo [\Mrd]}{Körperschaftsteuer: Saldo}{pos}
\kurve{kst}{bip}{Körperschaftsteuer [\%]}{BIP [\Mrd]}{Körperschaftsteuer: BIP der Periode}{neg}
\kurve{rentenniveau}{saldo}{Rentenniveau [\%]}{Saldo [\Mrd]}{Rentenniveau: Saldo}{pos}
\kurve{rentenniveau}{gini}{Rentenniveau [\%]}{Gini [Pkt.]}{Rentenniveau: Gini}{neg}
\kurve{co2}{saldo}{\COz-Preis [\euro/t]}{Saldo [\Mrd]}{\COz-Preis: Saldo}{pos}
\kurve{co2}{emiss}{\COz-Preis [\euro/t]}{Emissionen [Mt/a]}{\COz-Preis: Emissionen}{neg}
\end{groupplot}
\end{tikzpicture}
\caption{Saldo und eine zweite Kennzahl der ersten Periode über den ganzen Bereich des
Reglers, alle übrigen Regler im Status quo. Gestrichelt: Wert des Status quo. 21 Stützstellen
je Kurve.}
\label{fig:kurven}
\end{figure}

\paragraph{So liest man die Kurven.} Eine Gerade heißt: Jeder Schritt wirkt gleich, die
Gewichte oben gelten im ganzen Bereich. Ein Knick heißt: Dort greift eine Schwelle oder
Klammer der Engine (Abschnitt~\ref{sec:eigenschaften}). Eine flacher werdende Kurve heißt:
Die Verhaltensreaktion frisst einen wachsenden Teil der mechanischen Wirkung. Beispiele:
\begin{itemize}
\item \textbf{Spitzensteuersatz:} Unterhalb von 42\,\% wird die Proportionalzone auf den
  Spitzensatz gedeckelt, der Satz trifft nun die ganze obere Mitte, und der Saldo fällt
  steil. Oberhalb wirkt er nur im Pareto-Rand des obersten Prozents, und die Kurve ist flach
  und gerade. Die Ausweich- und Wegzugsreaktion des obersten Prozents setzt im ganzen
  Bereich \emph{nicht} ein (Abschnitt~\ref{sec:grenzen}).
\item \textbf{Umsatzsteuer:} leicht konkav, weil die Konsumreaktion mit dem Satz wächst und
  die Nachfrage mitzieht; eine Obergrenze erreicht sie im Bereich nicht.
\item \textbf{Körperschaftsteuer:} Das Aufkommen ist Satz mal schrumpfende Basis, also eine
  nach unten geöffnete Parabel; ihr Scheitel liegt weit außerhalb des Reglerbereichs, die
  Klammern der Investitionsreaktion werden nicht erreicht.
\item \textbf{\COz-Preis:} Die Emissionen fallen linear zwischen etwa 22\,\euro/t (darunter
  greift die obere Klammer) und etwa 255\,\euro/t (darüber greift die untere Klammer bei
  40\,\% der bepreisten Emissionen). Der Saldo zeigt eine \emph{Laffer-Kurve}: Das Aufkommen
  $p\cdot E(p)$ erreicht sein Maximum bei etwa 190\,\euro/t, weil die Bemessungsgrundlage
  schneller schrumpft, als der Preis steigt; jenseits der unteren Klammer steigt es wieder.
\end{itemize}

\section{Wirkung über den Pfad}

Tabelle~\ref{tab:pfad} zeigt, was nach fünf Perioden mit unveränderter Politik bleibt. Hier
wirken die Kanäle, die in der ersten Periode unsichtbar sind: das private Kapital, das
Arbeitsangebot im Niveau und vor allem der Zinseszins auf die Schulden.

\begin{table}[!htb]
\centering\small
% Erzeugt von docs/modell/messung.js — nicht von Hand bearbeiten
\begin{tabular}{lrrrrr}
\toprule
Stellgröße (Schritt) & \rotatebox{75}{Schuldenquote 2041 [Pp.]} & \rotatebox{75}{BIP 2041 [Mrd. \euro{}]} & \rotatebox{75}{CO\textsubscript{2} 2025–2040 [Mt]} & \rotatebox{75}{Saldo 2041–44 [Mrd. \euro{}/a]} & \rotatebox{75}{Gini 2041–44 [Gini-Pkt.]} \\
\midrule
Grundfreibetrag (1.000 \euro{}) & $+$0,85 & $+$28,2 & 0 & $-$3,6 & $-$0,124 \\
Eingangssteuersatz (1 Pp.) & $-$0,51 & $-$7,9 & 0 & $+$2,7 & $-$0,037 \\
Spitzensteuersatz (1 Pp.) & $-$0,09 & $-$0,2 & 0 & $+$0,6 & $-$0,018 \\
Umsatzsteuer (1 Pp.) & $-$2,44 & $-$8,2 & 0 & $+$14,7 & $+$0,057 \\
Erbschaftsteuer (1 Pp.) & $-$0,08 & $-$0,2 & 0 & $+$0,5 & $-$0,007 \\
Körperschaftsteuer (1 Pp.) & $-$0,33 & $-$8,1 & 0 & $+$1,4 & $-$0,038 \\
Gewerbesteuer (1 Pp.) & $-$0,93 & $-$9,6 & 0 & $+$5,0 & $-$0,072 \\
RV-Beitrag (1 Pp.) & $-$3,76 & $-$15,6 & 0 & $+$22,8 & $+$0,066 \\
KV-Beitrag (1 Pp.) & $-$3,90 & $-$14,5 & 0 & $+$23,6 & $+$0,122 \\
Rentenniveau (1 Pp.) & $+$2,01 & $+$9,0 & 0 & $-$13,3 & $-$0,275 \\
Bürgergeld (10 \euro{}/Monat) & $+$0,11 & $+$0,6 & 0 & $-$0,7 & $-$0,043 \\
Kindergeld (10 \euro{}/Monat) & $+$0,52 & $+$2,0 & 0 & $-$3,1 & $-$0,067 \\
CO\textsubscript{2}-Preis (10 \euro{}/t) & $-$0,46 & $-$0,8 & $-$103 & $+$1,7 & $+$0,007 \\
Klimageld (ein) & $+$2,19 & $+$4,3 & 0 & $-$8,2 & $-$0,119 \\
\bottomrule
\end{tabular}

\caption{Wirkung nach fünf Perioden (Ende des Pfads) gegenüber dem Status-quo-Pfad, je
Schritt, Politik über alle Perioden unverändert. Status quo: Schuldenquote \SQschuldEnde\,\%,
BIP \SQbipEnde\,\Mrd, \COz{} kumuliert \SQcoEnde\,Mt.}
\label{tab:pfad}
\end{table}

\paragraph{So liest man die Tabelle.} Teilen Sie die Änderung der Schuldenquote durch die
Saldowirkung der ersten Periode. Wirkt eine Stellgröße nur über den Saldo, liegt dieses
Verhältnis bei rund $-0{,}3$\,\Pp{} je Mrd.\,\euro{} (etwa Rentenbeitrag, Umsatzsteuer,
Kindergeld): Ein dauerhafter Saldo geht über 16 Jahre fast linear in die Quote ein. Weicht
das Verhältnis ab, ist ein weiterer Kanal im Spiel, der das BIP im Nenner bewegt:
\begin{itemize}
\item Der \textbf{Grundfreibetrag} kostet \wert{p1}{freibetrag}{saldo}\,\Mrd{} im Jahr, erhöht
  die Quote 2041 aber nur um \wert{pfad}{freibetrag}{schuld}\,\Pp: Das zusätzliche
  Arbeitsangebot hebt das BIP 2041 um \wert{pfad}{freibetrag}{bip}\,\Mrd.
\item Die \textbf{Körperschaftsteuer} bringt \wert{p1}{kst}{saldo}\,\Mrd, senkt die Quote aber
  nur um \wert{pfad}{kst}{schuld}\,\Pp: Das private Kapital schrumpft, das BIP 2041 sinkt um
  \wert{pfad}{kst}{bip}\,\Mrd.
\end{itemize}

\section{Die Messwerte im Einzelnen}

Zum Nachschlagen die vollständigen Messwerte für die Stellgrößen am Tisch.

\begin{table}[!htb]
\centering\scriptsize
\setlength{\tabcolsep}{3pt}
% Erzeugt von docs/modell/messung.js — nicht von Hand bearbeiten
\begin{tabular}{lrrrrrrrrrrr}
\toprule
Stellgröße (Schritt) & \rotatebox{75}{Einkommensteuer [Mrd. \euro{}]} & \rotatebox{75}{Umsatzsteuer [Mrd. \euro{}]} & \rotatebox{75}{KSt + GewSt [Mrd. \euro{}]} & \rotatebox{75}{Erb-, Boden-, Verm.-St. [Mrd. \euro{}]} & \rotatebox{75}{SV-Beiträge [Mrd. \euro{}]} & \rotatebox{75}{CO\textsubscript{2}-Einnahmen netto [Mrd. \euro{}]} & \rotatebox{75}{Transfers [Mrd. \euro{}]} & \rotatebox{75}{Rentenzahlungen [Mrd. \euro{}]} & \rotatebox{75}{Arbeitsangebot [\%]} & \rotatebox{75}{Investitionsfaktor [\%]} & \rotatebox{75}{Nachfrageimpuls [Mrd. \euro{}]} \\
\midrule
Grundfreibetrag (1.000 \euro{}) & $-$12,6 & $+$3,4 & $+$0,2 & $+$0,1 & $+$1,4 & 0,0 & 0,0 & 0,0 & $+$0,35 & 0,00 & $+$8,5 \\
Eingangssteuersatz (1 Pp.) & $+$4,5 & $-$1,1 & $-$0,1 & 0,0 & $-$0,5 & 0,0 & 0,0 & 0,0 & $-$0,08 & 0,00 & $-$2,8 \\
Spitzensteuersatz (1 Pp.) & $+$0,3 & 0,0 & 0,0 & 0,0 & 0,0 & 0,0 & 0,0 & 0,0 & 0,00 & 0,00 & $-$0,1 \\
Umsatzsteuer (1 Pp.) & $-$0,5 & $+$9,5 & $-$0,1 & 0,0 & $-$0,9 & 0,0 & 0,0 & 0,0 & 0,00 & 0,00 & $-$5,5 \\
Erbschaftsteuer (1 Pp.) & 0,0 & 0,0 & 0,0 & $+$0,3 & 0,0 & 0,0 & 0,0 & 0,0 & 0,00 & 0,00 & $-$0,1 \\
Körperschaftsteuer (1 Pp.) & $-$0,1 & $-$0,1 & $+$2,5 & 0,0 & $-$0,2 & 0,0 & 0,0 & 0,0 & 0,00 & $-$0,40 & $-$1,0 \\
Gewerbesteuer (1 Pp.) & $-$0,2 & $-$0,1 & $+$4,8 & 0,0 & $-$0,3 & 0,0 & 0,0 & 0,0 & 0,00 & $-$0,40 & $-$2,0 \\
RV-Beitrag (1 Pp.) & $-$0,9 & $-$2,3 & $-$0,3 & $-$0,1 & $+$15,7 & 0,0 & 0,0 & 0,0 & 0,00 & 0,00 & $-$10,3 \\
KV-Beitrag (1 Pp.) & $-$0,8 & $-$2,1 & $-$0,3 & $-$0,1 & $+$15,9 & 0,0 & 0,0 & 0,0 & 0,00 & 0,00 & $-$9,6 \\
Rentenniveau (1 Pp.) & $+$0,4 & $+$0,3 & $+$0,1 & 0,0 & $+$0,8 & 0,0 & 0,0 & $+$7,5 & 0,00 & 0,00 & $+$4,8 \\
Bürgergeld (10 \euro{}/Monat) & 0,0 & 0,0 & 0,0 & 0,0 & $+$0,1 & 0,0 & $+$0,5 & 0,0 & 0,00 & 0,00 & $+$0,4 \\
Kindergeld (10 \euro{}/Monat) & $+$0,1 & $+$0,1 & 0,0 & 0,0 & $+$0,2 & 0,0 & $+$2,0 & 0,0 & 0,00 & 0,00 & $+$1,3 \\
CO\textsubscript{2}-Preis (10 \euro{}/t) & $-$0,1 & $-$0,1 & 0,0 & 0,0 & $-$0,2 & $+$2,6 & 0,0 & 0,0 & 0,00 & 0,00 & $-$1,4 \\
Klimageld (ein) & $+$0,6 & $+$0,5 & $+$0,2 & $+$0,1 & $+$1,3 & $-$12,6 & $+$12,6 & 0,0 & 0,00 & 0,00 & $+$7,8 \\
\bottomrule
\end{tabular}

\caption{Wirkung auf die Kanäle (erste Periode, je Schritt). Das sind die Kantengewichte der
Ressortgraphen, einschließlich der dort weggelassenen kleinen Werte.}
\label{tab:kanaele}
\end{table}

\begin{table}[!htb]
\centering\scriptsize
\setlength{\tabcolsep}{3pt}
% Erzeugt von docs/modell/messung.js — nicht von Hand bearbeiten
\begin{tabular}{lrrrrrrrr}
\toprule
Stellgröße (Schritt) & \rotatebox{75}{Saldo [Mrd. \euro{}/a]} & \rotatebox{75}{Gini [Gini-Pkt.]} & \rotatebox{75}{Armutsrisiko [\%]} & \rotatebox{75}{Emissionen [Mt/a]} & \rotatebox{75}{BIP [Mrd. \euro{}]} & \rotatebox{75}{Netto D1 [\euro{}/a]} & \rotatebox{75}{Netto D5 [\euro{}/a]} & \rotatebox{75}{Netto D10c [\euro{}/a]} \\
\midrule
Grundfreibetrag (1.000 \euro{}) & $-$6,9 & $-$0,125 & $+$0,29 & 0,0 & $+$8,5 & 0 & $+$358 & $+$637 \\
Eingangssteuersatz (1 Pp.) & $+$2,6 & $-$0,036 & $-$0,09 & 0,0 & $-$2,8 & 0 & $-$101 & $-$322 \\
Spitzensteuersatz (1 Pp.) & $+$0,3 & $-$0,017 & 0,00 & 0,0 & $-$0,1 & 0 & 0 & $-$1.156 \\
Umsatzsteuer (1 Pp.) & $+$7,7 & $+$0,057 & $-$0,03 & 0,0 & $-$5,5 & $-$87 & $-$204 & $-$1.555 \\
Erbschaftsteuer (1 Pp.) & $+$0,3 & $-$0,007 & 0,00 & 0,0 & $-$0,1 & 0 & $-$2 & $-$177 \\
Körperschaftsteuer (1 Pp.) & $+$2,0 & $-$0,035 & $-$0,01 & 0,0 & $-$1,0 & $-$8 & $-$27 & $-$1.616 \\
Gewerbesteuer (1 Pp.) & $+$3,9 & $-$0,068 & $-$0,02 & 0,0 & $-$2,0 & $-$16 & $-$52 & $-$3.128 \\
RV-Beitrag (1 Pp.) & $+$11,7 & $+$0,064 & $-$0,10 & 0,0 & $-$10,3 & $-$129 & $-$366 & $-$987 \\
KV-Beitrag (1 Pp.) & $+$12,1 & $+$0,119 & $-$0,10 & 0,0 & $-$9,6 & $-$129 & $-$366 & $-$679 \\
Rentenniveau (1 Pp.) & $-$5,8 & $-$0,216 & $-$0,01 & 0,0 & $+$4,8 & $+$122 & $+$238 & $+$117 \\
Bürgergeld (10 \euro{}/Monat) & $-$0,4 & $-$0,042 & $-$0,08 & 0,0 & $+$0,4 & $+$72 & 0 & 0 \\
Kindergeld (10 \euro{}/Monat) & $-$1,7 & $-$0,066 & $-$0,02 & 0,0 & $+$1,3 & $+$45 & $+$59 & $+$11 \\
CO\textsubscript{2}-Preis (10 \euro{}/t) & $+$2,0 & $+$0,019 & $-$0,01 & $-$9,8 & $-$1,4 & $-$24 & $-$55 & $-$303 \\
Klimageld (ein) & $-$9,7 & $-$0,317 & $-$0,15 & 0,0 & $+$7,8 & $+$307 & $+$307 & $+$307 \\
\bottomrule
\end{tabular}

\caption{Wirkung auf die Kennzahlen (erste Periode, je Schritt).}
\label{tab:kennzahlen}
\end{table}

\begin{table}[!htb]
\centering\scriptsize
\setlength{\tabcolsep}{3pt}
% Erzeugt von docs/modell/messung.js — nicht von Hand bearbeiten
\begin{tabular}{lrrrrrrrr}
\toprule
Stellgröße (Schritt) & \rotatebox{75}{Saldo [Mrd. \euro{}/a]} & \rotatebox{75}{Gini [Gini-Pkt.]} & \rotatebox{75}{Armutsrisiko [\%]} & \rotatebox{75}{Emissionen [Mt/a]} & \rotatebox{75}{BIP [Mrd. \euro{}]} & \rotatebox{75}{Netto D1 [\euro{}/a]} & \rotatebox{75}{Netto D5 [\euro{}/a]} & \rotatebox{75}{Netto D10c [\euro{}/a]} \\
\midrule
Grenze Spitzensatz ($-$10.000 \euro{}) & 0,0 & $-$0,001 & 0,00 & 0,0 & 0,0 & 0 & 0 & $-$92 \\
USt ermäßigt (1 Pp.) & $+$4,4 & $+$0,030 & $-$0,02 & 0,0 & $-$2,9 & $-$46 & $-$108 & $-$824 \\
Abgeltungsteuer (1 Pp.) & $+$2,3 & $-$0,072 & $-$0,01 & 0,0 & $-$0,8 & $-$1 & $-$11 & $-$2.984 \\
Beitragsbemessungsgr. (10.000 \euro{}) & $+$6,0 & $-$0,274 & 0,00 & 0,0 & $-$3,6 & 0 & 0 & $-$1.704 \\
Vermögensteuer (0,1 Pp.) & $+$2,8 & $-$0,075 & $-$0,02 & 0,0 & $-$1,3 & 0 & $-$19 & $-$1.868 \\
Bodenwertsteuer (0,1 Pp.) & $+$4,2 & $-$0,109 & $-$0,02 & 0,0 & $-$2,0 & 0 & $-$27 & $-$2.726 \\
Zucman-Mindeststeuer (1 Pp.) & $+$3,4 & $-$0,142 & 0,00 & 0,0 & $-$0,7 & 0 & 0 & $-$9.476 \\
\bottomrule
\end{tabular}

\caption{Wirkung der Stellgrößen der klassischen Fläche auf die Kennzahlen (erste Periode).}
\label{tab:weitere}
\end{table}


\clearpage
\appendix
\part*{Anhang}
\addcontentsline{toc}{part}{Anhang}
% SPDX-License-Identifier: CC-BY-4.0
% Copyright 2025 Florian Aram Feuerriegel — kassensturz.org
% Anhang — Daten, Elastizitäten, Quellen (alle Tabellen erzeugt von messung.js)

\section{Haushaltsdaten}
\begin{table}[!htb]
\centering\footnotesize
\setlength{\tabcolsep}{3.5pt}
% Erzeugt von docs/modell/messung.js — nicht von Hand bearbeiten
\begin{tabular}{lrrrrrrrrrrr}
\toprule
Dezil & $B_i$ & $\kappa_i$ & $c_i$ & $w_i$ & $V_i$ & $N_i$ & $\rho_i$ & $m_i$ & $k_i$ & $q_i$ & $\ell_i$ \\
 & Tsd.\,\euro{} & \% & \% & & Tsd.\,\euro{} & Mio. & \% & & & & \euro{} \\
\midrule
D1 & 14 & 1 & 100 & 1,3 & 1 & 4,10 & 42,0 & 0,65 & 0,37 & 0,60 & 166 \\
D2 & 21 & 1 & 99 & 1,4 & 5 & 4,10 & 46,2 & 0,62 & 0,51 & 0,25 & 236 \\
D3 & 27 & 2 & 96 & 1,5 & 15 & 4,10 & 42,0 & 0,58 & 0,56 & 0,08 & 287 \\
D4 & 33 & 2 & 92 & 1,6 & 35 & 4,10 & 35,3 & 0,55 & 0,54 & 0,02 & 332 \\
D5 & 40 & 3 & 88 & 1,7 & 70 & 4,10 & 28,6 & 0,52 & 0,49 & 0,00 & 379 \\
D6 & 48 & 3 & 85 & 1,8 & 120 & 4,10 & 21,8 & 0,48 & 0,44 & 0,00 & 426 \\
D7 & 58 & 4 & 82 & 1,9 & 200 & 4,10 & 15,1 & 0,45 & 0,40 & 0,00 & 480 \\
D8 & 72 & 5 & 78 & 2,0 & 340 & 4,10 & 10,1 & 0,42 & 0,35 & 0,00 & 532 \\
D9 & 95 & 7 & 72 & 2,0 & 620 & 4,10 & 6,7 & 0,38 & 0,30 & 0,00 & 618 \\
D10a & 125 & 8 & 60 & 2,0 & 850 & 2,05 & 4,2 & 0,32 & 0,23 & 0,00 & 665 \\
D10b & 220 & 18 & 52 & 2,0 & 1.500 & 1,64 & 2,5 & 0,25 & 0,16 & 0,00 & 976 \\
D10c & 700 & 45 & 40 & 2,0 & 7.000 & 0,41 & 0,8 & 0,15 & 0,09 & 0,00 & 2.070 \\
\bottomrule
\end{tabular}

\caption{Die zwölf Haushaltstypen (\datei{DEZILE} und Profile in \datei{data.js}).
$B_i$ Bruttoeinkommen, $\kappa_i$ Kapitalanteil, $c_i$ durchschnittliche Konsumquote,
$w_i$ Bedarfsgewicht (neue OECD-Skala), $V_i$ Vermögen, $N_i$ Haushalte,
$\rho_i$ Rentenanteil am Einkommen (Näherung), $m_i$ Grenzkonsumneigung (Näherung),
$k_i$ Kindergeldkinder je Haushalt, $q_i$ Bürgergeld-Regelsatzjahre je Haushalt,
$\ell_i$ \COz-Last je Haushalt im Status quo (55\,\euro{}/t, vollständig überwälzt).}
\label{tab:dezile}
\end{table}

\section{Elastizitäten}
{\small\sloppy% Erzeugt von docs/modell/messung.js — nicht von Hand bearbeiten
\begin{longtable}{p{2.6cm}rp{2.2cm}p{8.2cm}}
\toprule
Schlüssel & Wert & Spanne & Fundstelle \\
\midrule
\endhead
\texttt{labor\_supply} & 0,200 & 0,1–0,3 & Saez/Chetty/Gruber Konsens · ifo Schnelldienst 01/2025 \\
\texttt{capital\_supply} & 0,500 & 0,4–0,8 & Kleven/Schultz (2014) JPubEc \\
\texttt{consumption} & $-$0,350 & −0,2 bis −0,5 & Lewbel/Pendakur (2009) JPubEc · Metaanalyse Havranek et al. 2018 \\
\texttt{co2} & $-$0,300 & −0,2 bis −0,4 & EWI/DIW BEHG-Evaluation 2023 · Edenhofer/PIK 2024 \\
\texttt{evasion} & 0,250 & 0,1–0,3 & Schneider (2023) Shadow Economy DE · IfW Kiel 2024 \\
\texttt{investment} & $-$0,400 & −0,3 bis −0,5 & Gechert/Heimberger (2022) NIER · Neumeier SVR Arbeitspapier 03/2025 \\
\texttt{d10c\_labor} & 0,400 & 0,25–0,5 & Saez/Slemrod/Giertz (2012) JEL 50(1) · Piketty/Saez/Stantcheva (2014) AEJ: Economic Policy 6(1) \\
\texttt{d10c\_avoidance} & 0,500 & 0,3–0,7 & Kleven/Schultz (2014) JPubEc · Chetty/Friedman/Saez (2013) \\
\texttt{d10c\_wegzug} & 0,100 & 0,05–0,15 & Beleg ausstehend \\
\texttt{verm\_ausweichen} & 0,215 & 0,1–0,43 je Pp. & Brülhart/Gruber/Krapf/Schmidheiny (2022) AEJ: Economic Policy 14(4), Behavioral Responses to Wealth Taxes \\
\bottomrule
\end{longtable}
}

\section{Formelquellen}
Die Objekte \datei{FORMEL\_QUELLEN\_*} der Engine, unverändert abgedruckt.
{\small\sloppy% Erzeugt von docs/modell/messung.js — nicht von Hand bearbeiten
\begin{longtable}{p{3.2cm}p{11.6cm}}
\toprule
Quelle & Formel und Fundstelle \\
\midrule
\endhead
\multicolumn{2}{l}{\textbf{\texttt{berechne.js}}} \\
\texttt{arbeitsangebot} & {\formelschrift labor\_factor = 1 + ε × Δ(1 − GS) / (1 − GS\_SQ)}\newline Saez/Chetty/Gruber Konsens · ifo Schnelldienst 01/2025 · Piketty/Saez/Stantcheva (2014) AER \\
\texttt{bge\_arbeitsangebot} & {\formelschrift lf = lf\_tax × (1 − BGE\_LABOR\_EFF[i] × min(1,67; BGE/1200))}\newline RWI (2024) bis −30 \% bei 1.500 \euro{} · DIW Pilot 2024 (n=107) · ZEW Heim et al. (extensiver Margin) · ifo Mikrosimulation 2021 \\
\texttt{mwst\_konsumanteil} & {\formelschrift MwSt = Konsum × (0,70 × t\_reg/(1+t\_reg) + 0,30 × t\_erm/(1+t\_erm))}\newline Destatis VGR 2024 (ca. 70 \% Regelsatz-Konsum) · Lewbel/Pendakur (2009) JPubEc \\
\texttt{co2\_emissionen} & {\formelschrift Emissionen = Pfad(Jahr) × [(1 − s) + s × max(0,4; min(1,1; 1 + ε\_CO₂ × (p − 55)/100))],  s = 327/649}\newline BEHG § 10 · EWI/DIW BEHG-Evaluation 2023 · Edenhofer/PIK 2024 · UBA Projektionsbericht 2025 (Pfad) · UBA Treibhausgas-Emissionen 2025 \\
\texttt{erbschaft} & {\formelschrift Erb\_auf = Masse\_top × Satz\_eff + Masse\_unten × min(Satz,15\%)/2}\newline Bach/Sinclair/Bührle/Wichers DIW 2025 · §§ 10, 19 ErbStG · BRH 2016 \\
\texttt{zucman} & {\formelschrift Zucman\_auf = Milliardärsvermögen × max(0; Satz − 0,3 \% ESt − VermSt-Satz) × (1 − 0,15 × min(1; Satz/2))}\newline Zucman (2024) A blueprint for a coordinated minimum effective taxation standard for ultra-high-net-worth individuals, G20-Bericht · EU Tax Observatory 2024 · Jakobsen/Kleven/Kolsrud NBER 2024 \\
\texttt{sv\_beitraege} & {\formelschrift SV = Lohnsumme\_sv × Satz\%  (nur bis BBG)}\newline § 158 SGB VI · § 241 SGB V · § 341 SGB III · § 55 SGB XI · DRV Beitragssätze 2025 \\
\texttt{rv\_ausgaben} & {\formelschrift ΔRente\_i = Brutto\_i × rente\_anteil\_i × renten\_faktor × (Rentenniveau/48 \% − 1);  ΔRV-Ausgaben = Σ Haushalte ΔRente\_i;  KV, AL, PV fest}\newline § 68 SGB VI (Rentenanpassung) · § 154 Abs. 3 SGB VI (Sicherungsniveau) · Rentenpaket 2025 (BMAS, Haltelinie 48 \% bis 2031) · § 158 SGB VI · BMAS Rentenversicherungsbericht 2025 \\
\texttt{verwaltungskosten} & {\formelschrift Admin = ∑ Aufkommen\_i × Quote\_i  (ADMIN\_QUOTE je Steuer)}\newline BRH Jahresberichte · OECD Tax Administration 2024 · Normenkontrollrat Jahresbericht 2024 \\
\texttt{deadweight\_loss} & {\formelschrift DWL = ∑ᵢ 0,5 × ε × GS\_i² / (1 − GS\_i) × Lohnsumme\_i  (Harberger-Dreieck je Dezil)}\newline Harberger (1964) · Saez (2001) JPubEc · Chetty (2009) AER P\&P \\
\texttt{dynamisches\_scoring} & {\formelschrift Δ\_dyn = Δ\_KSt × (investment\_factor − 1) + Δ\_ESt × (avg\_labor − 1)}\newline CBO Dynamic Scoring Guidelines · ifo Schnelldienst 01/2025 · Saez/Chetty Konsens \\
\texttt{inzidenz} & {\formelschrift Last\_i = ΔKSt+ΔGewSt × (½ Anteil Arbeitseinkommen\_i + ½ Anteil Kapitaleinkommen\_i) + ΔErbSt+ΔVermSt+ΔBodenSt × Anteil Vermögen\_i + ΔZucman × D10c}\newline Fuest/Peichl/Siegloch (2018) AER 108(2): Beschäftigte tragen rund 51 \% der Gewerbesteuer · Harberger (1962) · CBO (2012) Distribution of Corporate Tax · Mirrlees Review (2011) Tax by Design, Kap. 16 (Bodenwertsteuer) \\
\midrule
\multicolumn{2}{l}{\textbf{\texttt{einkommensteuer.js}}} \\
\texttt{estTarif} & {\formelschrift ∫₀ˣ r(z) dz  (stückweise lineare Grenzsteuerrate, 5 Zonen)}\newline § 32a Abs. 1 EStG 2026 (Steuerfortentwicklungsgesetz, BGBl. 2024 I Nr. 449) · Formeltarif (kontinuierlich, keine Sprünge) \\
\texttt{grenzsteuersatz} & {\formelschrift r(z) = r₀ + (rₘ − r₀) × z/g₂  [Zone 2], linear interpoliert je Zone}\newline § 32a Abs. 1 Nr. 2–5 EStG · BMF Steuerschätzung 2025 \\
\texttt{effSteuersatz} & {\formelschrift T(z) / z}\newline § 2 Abs. 5 EStG \\
\texttt{spitzenzone} & {\formelschrift Zuschlag\_D10c = (r5 − r4) × E[max(0, x − Grenze)],  x \textasciitilde{} Pareto(α = 1,5) mit dem zvE-Mittel des Dezils}\newline Atkinson/Piketty/Saez (2011) JEL 49(1), Top Incomes in the Long Run of History · Bartels/Jenderny (2015) DIW Discussion Paper 1508 \\
\texttt{estHaushalt} & {\formelschrift T\_HH = s × T(brutto × q / s)  mit q = zvE-Quote, s = Splitting-Faktor}\newline §§ 26, 32a Abs. 5 EStG (Splitting) · Destatis ESt-Statistik (zvE vs. Bruttoeinkünfte) · Mikrozensus 2024 \\
\midrule
\multicolumn{2}{l}{\textbf{\texttt{verteilung.js}}} \\
\texttt{aequivalenzEinkommen} & {\formelschrift y\_äq,i = Netto\_i / Bedarfsgewicht\_i  (neue OECD-Skala: 1,0 erste Person · 0,5 weitere ab 14 J. · 0,3 Kinder unter 14 J.)}\newline Eurostat EU-SILC Methodik (äquivalisiertes verfügbares Einkommen) · Destatis Glossar „Äquivalenzeinkommen" · OECD (2013) Framework for Statistics on the Distribution of Household Income \\
\texttt{berechneGini} & {\formelschrift G = 1 − 2·∫Lorenz(x)dx  (Trapezregel, gewichtet nach Haushaltszahl)}\newline Sen (1973) On Economic Inequality · Cowell (2011) Measuring Inequality · Destatis Methodik Gini-Koeffizient \\
\texttt{berechnePalma} & {\formelschrift Palma = Einkommensanteil(Top 10\%) / Einkommensanteil(Bottom 40\%)}\newline Palma (2011) Homogeneous Middles vs. Heterogeneous Tails · UNDP HDR 2013 \\
\texttt{berechneMedianGewichtet} & {\formelschrift Gewichteter Median: kumulierte Haushaltsanteile bis 50 \%}\newline Destatis Mikrozensus 2024 · SOEP v40 \\
\texttt{armutsrisiko} & {\formelschrift pov\_i = pov\_sq\_i × (y\_i/y\_sq\_i ÷ PL/PL\_sq)\textasciicircum{}(−1,5)}\newline Bourguignon (2003) · EU-SILC DE 2023 (14,8 \% Kalibrierung) · SOEP v40 · IAB Kurzbericht 2024 \\
\texttt{berechneNettoSQ} & {\formelschrift Netto\_SQ = Brutto − ESt(SQ) − SV(SQ) − MwSt(SQ) − CO₂(SQ) + Klimageld(SQ) + Transfers(SQ)}\newline PRESETS.status\_quo (data.js) · § 32a EStG 2025 · § 158 SGB VI · § 241 SGB V \\
\texttt{berechneDezilDelta} & {\formelschrift Δ\_i = Netto\_neu\_i − Netto\_SQ\_i}\newline SOEP v40 · Destatis Mikrozensus 2024 · BMF Steuerschätzung 2025 \\
\bottomrule
\end{longtable}
}


\end{document}
